\documentclass[10pt,a4paper]{article}
\usepackage{amsmath,amssymb,amsthm,graphicx}
\usepackage{amsmath,amssymb,amsthm,hyperref,xfrac}
\providecommand{\skpt}{\smallskip}
\providecommand{\ie}{i.e.\ }
\providecommand{\eg}{e.g.\ }
\providecommand{\resp}{respectively\ }
\providecommand{\Subsection}{\subsection}
\providecommand{\Subsubsection}{\subsubsection}
\usepackage[scr=rsfs,cal=euler]{mathalfa}
\multlinegap0pt
\renewcommand{\theenumi}{\alph{enumi}}
\renewcommand{\theenumii}{\roman{enumii}}
\newenvironment{enumeratei}
{\bgroup\def\theenumi{\roman{enumi}}\def\theenumii{\arabic{enumii}}\begin{enumerate}}
{\end{enumerate}\egroup}
\makeatletter
\let\p@enumii\relax
\makeatother
\newcommand{\ldot}{\mathbin{.}}
\newcommand{\lcdots}{\ldot\cdots\ldot}
\newcommand{\RedefinitSymbole}[1]{%
\expandafter\let\csname old\string#1\endcsname=#1
\let#1=\relax
\newcommand{#1}{\csname old\string#1\endcsname\,}%
}
\RedefinitSymbole{\forall} \RedefinitSymbole{\exists}

\newcommand\mto{\mathchoice{\longmapsto}{\mapsto}{\mapsto}{\mapsto}}
\newcommand{\psfrac}[2]{\sfrac{(#1)}{#2}}
\newcommand{\spfrac}[2]{\sfrac{#1}{(#2)}}

\usepackage[all,cmtip]{xy}\let\labelstyle\textstyle
\usepackage{stmaryrd}
\usepackage{bbm}
\usepackage{mathtools}

\theoremstyle{plain}
\newtheorem{themBJ}{Theorem}
\renewcommand*{\thethemBJ}{\Alph{themBJ}}
\newtheorem{them}{Theorem}[section]
\newtheorem{pro}[them]{Proposition}
\newtheorem{pronotation}[them]{Proposition/Notation}
\newtheorem{prodefinition}[them]{Proposition/Definition}
\newtheorem{lemnotation}[them]{Lemma/Notation}
\newtheorem{cor}[them]{Corollary}
\newtheorem{lem}[them]{Lemma}
\newtheorem*{claim}{Claim}

\theoremstyle{definition}
\newtheorem{vocabulary}[them]{Vocabulary}
\newtheorem{defi}[them]{Definition}
\newtheorem{convention}[them]{Convention}
\newtheorem{notation}[them]{Notation} %?
\newtheorem{remnotation}[them]{Remark/Notation} %?
\newtheorem{definotation}[them]{Definition/Notation}
\newtheorem{remdefi}[them]{Remark/Definition}
\newtheorem{reminder-notation}[them]{Reminder/Notation}
\newtheorem{rem}[them]{Remark}
\newtheorem{ex}[them]{Example}
\newtheorem{reminder}[them]{Reminder}
\newtheorem*{problem*}{Main Problem}
\newtheorem*{note}{Note}
\newtheorem*{fact}{Fact}

\newcommand{\dd}{\mathrm{d}}
\newcommand{\rec}{\mathfrak R}
\newcommand{\id}{\mathrm{Id}}
\newcommand{\join}{\mathrm{Joint}}
\newcommand{\como}{\mathfrak{Q}}
\newcommand{\mq}{\mathfrak{MQ}}
\newcommand{\lev}{\mathrm{Lev}}
\newcommand{\ML}{\mathrm{LK}}
\newcommand{\IT}{\mathrm{IK}}
\newcommand{\disp}{\mathfrak{Disp}}
\newcommand{\argth}{\operatorname{Argth}}
\newcommand{\vol}{\operatorname{Vol}}
\newcommand{\parti}{\operatorname{Part}}
\newcommand{\idk}{\operatorname{id}}

\newcommand{\espace}{E}
\newcommand{\subd}{R}
\newcommand{\Tt}{{\mathcal T}}
\newcommand{\XX}{{\mathcal X}}
\newcommand{\EE}{{\mathcal E}}
\newcommand{\A}{{\mathcal A}}
\newcommand{\Var}{{\mathcal L}}
\newcommand{\Id}{\operatorname{Id}}
\newcommand{\one}{\mathbbm{1}}
\let\mb\relax
\let\mn\relax
\newcommand{\trans}{{}^{\mathrm{t}}\mkern-1mu}
\newcommand{\eqdef}{\coloneqq}
\newcommand{\eqdefup}{\coloneqq}
\newcommand{\espcdot}{\mbox{$\,\cdot\,$}}

\let\eps\varepsilon

\newcommand{\Lg}{\lambda|_{[0,1]}}
\newcommand{\la}{\lambda}
\renewcommand{\P}{\mathbb{P}}
\newcommand{\p}{\mathcal{P}}
\newcommand{\m}{\mathcal{M}}
\newcommand{\ma}{\mathrm{Marg}}
\newcommand{\cC}{\mathcal{C}}\let\C\cC

\newcommand{\R}{\mathbb{R}}
\newcommand{\Q}{\mathbb{Q}}
\newcommand{\N}{\mathbb{N}}
\newcommand{\Z}{\mathbb{Z}}

\newcommand{\AND}{\quad\text{and}\quad}

\newcommand{\E}{\mathbb{E}}
\newcommand{\law}{\operatorname{Law}}
\newcommand{\var}{\operatorname{Var}}
\newcommand{\sto}{\mathrm{sto}}
\newcommand{\stosup}{\operatorname{stosup}}
\newcommand{\stoinf}{\operatorname{stoinf}}
\newcommand{\lcsup}{\operatorname{losup}}
\newcommand{\lcinf}{\operatorname{loinf}}

\newcommand{\leqp}{\preceq_+}
\newcommand{\leqc}{\preceq_{C}}
\newcommand{\leqcs}{\preceq_{C,\sto}}
\newcommand{\leqs}{\preceq_{\text{\rm sto}}}
\newcommand{\geqs}{\succeq_{\text{\rm sto}}}
\newcommand{\e}{\mathrm{e}}

\newcommand{\pr}{\operatorname{proj}}

\newcommand{\pdec}{\mathcal{P}^{\scriptscriptstyle\searrow}}
\newcommand{\mdec}{\mathcal{M}^{\scriptscriptstyle\searrow}}
\newcommand{\leb}{\LL_{[0,1]}}
\newcommand{\op}{\mathopen}
\newcommand{\clo}{\mathclose}
\newcommand{\leqlc}{\preceq_{\rm lo}}
\newcommand{\geqlc}{\succeq_{\rm lo}}
\newcommand{\rsemibracket}{\mathclose{\rfloor\hspace*{-.9ex}\lceil}}
\newcommand{\lsemibracket}{\mathclose{\rfloor\hspace*{-.9ex}\lceil}}
\newcommand{\semi}{\mathclose{\rfloor\hspace*{-.8ex}\lceil}}

\newlength{\arraycolsepsauvegardegenerale}
\setlength{\arraycolsepsauvegardegenerale}{\arraycolsep}
\newlength{\arraycolsepsauvegardetemporaire}

\newlength{\tabcolsepsauvegardegenerale}
\setlength{\tabcolsepsauvegardegenerale}{\tabcolsep}
\newlength{\tabcolsepsauvegardetemporaire}

\newenvironment{nnarray}%
{\setlength{\arraycolsepsauvegardetemporaire}{\arraycolsep}\setlength{\arraycolsep}{0cm}\begin{array}}
{\end{array}\setlength{\arraycolsep}{\arraycolsepsauvegardetemporaire}}

\newenvironment{nntabular}%
{\setlength{\tabcolsepsauvegardetemporaire}{\tabcolsep}\setlength{\tabcolsep}{0cm}\begin{tabular}}
{\end{tabular}\setlength{\tabcolsep}{\tabcolsepsauvegardetemporaire}}
\usepackage{cleveref}
\crefname{originalTheorem}{original Theorem}{original Theorems}
\Crefname{originalTheorem}{Original Theorem}{Original Theorems}
\makeatletter
\newlabel{them:a}{{A}{source}{}{sourcepointer-1}{}}
\newlabel{them:a@cref}{{[originalTheorem][0][]A}{[1][0][]source}{}{}{}}
\makeatother
\crefname{originalsection}{original section}{original sections}
\Crefname{originalsection}{Original section}{Original sections}
\makeatletter
\newlabel{item:limite_de_compose}{{A(a)(iv)}{source}{}{sourcepointer-2}{}}
\newlabel{item:limite_de_compose@cref}{{[originalsection][0][]A(a)(iv)}{[1][0][]source}{}{}{}}
\makeatother
\crefname{originalsection}{original section}{original sections}
\Crefname{originalsection}{Original section}{Original sections}
\makeatletter
\newlabel{item:them:deuxieme_c}{{4.21(c)}{source}{}{sourcepointer-3}{}}
\newlabel{item:them:deuxieme_c@cref}{{[originalsection][0][]4.21(c)}{[1][0][]source}{}{}{}}
\makeatother
\crefname{originalProposition}{original Proposition}{original Propositions}
\Crefname{originalProposition}{Original Proposition}{Original Propositions}
\makeatletter
\newlabel{prodef:transitions}{{3.11}{source}{}{sourcepointer-4}{}}
\newlabel{prodef:transitions@cref}{{[originalProposition][0][]3.11}{[1][0][]source}{}{}{}}
\makeatother
\crefname{originalsection}{original section}{original sections}
\Crefname{originalsection}{Original section}{Original sections}
\makeatletter
\newlabel{prodef:transitionsa}{{3.11(a)}{source}{}{sourcepointer-5}{}}
\newlabel{prodef:transitionsa@cref}{{[originalsection][0][]3.11(a)}{[1][0][]source}{}{}{}}
\makeatother
\crefname{originalsection}{original section}{original sections}
\Crefname{originalsection}{Original section}{Original sections}
\makeatletter
\newlabel{prodef:transitionsb}{{3.11(b)}{source}{}{sourcepointer-6}{}}
\newlabel{prodef:transitionsb@cref}{{[originalsection][0][]3.11(b)}{[1][0][]source}{}{}{}}
\makeatother
\crefname{originalsection}{original section}{original sections}
\Crefname{originalsection}{Original section}{Original sections}
\makeatletter
\newlabel{rem:qualitative}{{4.8}{source}{}{sourcepointer-7}{}}
\newlabel{rem:qualitative@cref}{{[originalsection][0][]4.8}{[1][0][]source}{}{}{}}
\makeatother
\crefname{originalsection}{original section}{original sections}
\Crefname{originalsection}{Original section}{Original sections}
\makeatletter
\newlabel{rem:stosup_existe}{{3.9}{source}{}{sourcepointer-8}{}}
\newlabel{rem:stosup_existe@cref}{{[originalsection][0][]3.9}{[1][0][]source}{}{}{}}
\makeatother
\crefname{originalsection}{original section}{original sections}
\Crefname{originalsection}{Original section}{Original sections}
\makeatletter
\newlabel{subsec:stochastic_order}{{3.1}{source}{}{sourcepointer-9}{}}
\newlabel{subsec:stochastic_order@cref}{{[originalsection][0][]3.1}{[1][0][]source}{}{}{}}
\makeatother
\title{A single nested sequence of finite Markovization times approximates the Markov-quantile process}
\author{Charles Boubel, Nicolas Juillet}
\date{doi:10.5802/jep.177; CC BY 4.0}
\providecommand{\oldbigcup}{\bigcup}
\begin{document}
\maketitle
\noindent\textit{Editorial note.} This task counts Theorem B in the author stated finite-dimensional convergence sense. No arbitrary dense time set, full path topology, continuous version or strong Markov assertion is claimed. The detailed theorem excerpt stops at(a),(b); its(c) classification is not assumed.\par\medskip
\noindent\textit{Editorial note.} Increasing kernel is used in Proposition3.11(b), printed here. The unused probability-valued monotone extension to every point outside the initial support in3.11(c) is not asserted. All actual level kernels act on[0,1]; general quantile images use almost-sure versions.\par\medskip
\noindent\textit{Editorial note.} The source bounded-density closure wording1095-1104 is used only with a common bound. In1345 a fixed input theta<=M*mu has theta P\_n<=M*nu because the target marginal is fixed; thus the required uniform bound is present. No closure claim for arbitrarily unbounded densities is assumed.\par\medskip
\noindent\textit{Editorial note.} At970 the general d-dimensional rectangle increment has the usual sign(-1)\textasciicircum{}(d-|epsilon|); the printed even-parity convention is correct in dimension2, the dimension of the directed level-coupling construction. At993 the mixed-increment condition must hold for the infimum CDF, as991 and the actual proof1017 explain, not merely for each individual CDF.\par\medskip
\noindent\textit{Editorial note.} At1605 the CDF comparison is reversed in print: the stochastic monotonicity proved1582 gives F[L\_R]<=F[L\_R-prime] when R-prime subset R, consistent with the stated lower-orthant order1591. Original formulas are retained.\par\medskip
\noindent\textit{Editorial note.} At1817 the finite approximating choices can and must retain the preceding finite set when moving to accuracy n, as required by1811. Equivalently take cumulative finite unions of the displayed choices; monotonicity1804/1591 preserves every approximation. This makes the author intended nested diagonal construction explicit.1876-1877 uses an undefined R-prime\_n; the directly preceding1811 approximation and1804 limit already identify the limit for the chosen R\_n. Do not infer a common time set from pairwise existence alone.\par\medskip
\noindent\textit{Editorial note.} At1832 the second factor concerns(t,s-prime), as correctly written in the next line1833, not the repeated(s,t) interval. The contradiction remains the author composition/continuity argument.\par\medskip
\noindent\textit{Editorial note.} In the atomic covariance proof,1915 suppresses the Lebesgue factor in the constant-density measure, and1919 reverses its displayed indicator inequality. Both are fixed by1915-1917 and the nonnegative covariance identity. At1957 both lower-orthant kernel test functions are decreasing, not increasing; the same covariance inequality holds for two functions with the same direction, consistently with2062-2063.\par\medskip
\noindent\textit{Editorial note.} The multivariate targets in1925/2054 have the stated coordinate marginals but are not assumed independent. The actual kernel integrals1941-1966 and the chain formula2063 use their full joint laws. The printed tensor-product target notation cannot be read as an extra product-law hypothesis.\par\medskip
\noindent\textit{Editorial note.} In1994-2016, m/n both refer to the number of Markovization times; all blocks B\_0,...,B\_m are included. The explicit finite endpoint products in2000-2002 determine the block law; the L\_[r\_k,r\_(k+1)] label there is not a replacement by all intervening real times or an identification of two- and three-coordinate laws. At2005 the defining q\_r has starting measure lambda, as1454-1456 specifies, not mu\_r.2016 ends its interior list atell\_k, as correctly written2021, and2051 ends the chain withs\_(k-1),s\_k.\par\medskip
\noindent\textit{Editorial note.} Original references to the simultaneous two-point/whole-process equality are not treated as proof shortcuts: the full finite-level pushforward argument1990-2036 and the atomic-factorization argument2050-2064 are retained.\par\medskip
\noindent\textit{Editorial note.} Source1375 says the sequence is increasing: the needed and definitionally supplied fact is that each function tilde(h)\_n is increasing in its spatial variable, not that n maps to that function is increasing.\par\medskip
\noindent\textit{Editorial note.} The pointwise proof1373-1385 uses [0,1] notation while the concatenation result is on R. Its bounded-coordinate argument transports through any fixed increasing homeomorphism R to(0,1); this ordinary coordinate convention and the already printed quantile transform handle general probability marginals. Independent review must confirm that no substantive missing author argument is concealed by this convention.\par\medskip
\noindent\textit{Editorial note.} At1381 the symmetric lower-limit inequality is printed <=; the symmetric claim needed for convergence is >=, uniquely fixed by the limsup bound being proved and the target conclusion. At1379 the displayed difference of half-line indicators gives[a,b), not[a,b]; either endpoint convention is allowed among monotone test functions, including1\_(b,infinity) for the closed version. The originals remain visible.\par\medskip
\noindent\textit{Editorial note.} At1389 the displayed order P-prime circ P reverses the typed concatenation. The defining marginals and actual proof1395,1399 unambiguously use P circ P-prime. At1402 tilde(h)\_n means P-prime\_n dot h, as the preceding integral1399 explicitly shows. At1435 the repeated index i+i and last n are local index slips determined by the adjacent projection i,i+1 and chain length d. No author formula is silently overwritten.\par\medskip
\noindent\textit{Editorial note.} At1807 the final displayed distance has the finite approximant as its first entry. The initial monotone-metric remark in the same proof directly bounds rho(L\_Rtau,L\_R) by rho(L\_Rprime,L\_R); this is the quantity required for convergence. The source order sandwich and proof are retained.\par\medskip
\noindent\textit{Editorial scope.} One Theorem B counting unit. The detailed Theorem4.21(a),(b) is supporting proof; its separate(c) classification/necessity statements and later qualitative examples are not additional tasks. The full source1451-1705 construction is retained to define the claimed limit without assuming the sought approximation.\par\medskip
\noindent\textit{Allowed standard prerequisites (editorial).}\begin{itemize}
\item Disintegration, Fubini, dominated convergence, monotone classes and Kolmogorov-Daniell extension. General probability foundations; the particular consistency of level kernels and preservation of the Markov property after quantile pushforward remain proved.
\item Portmanteau theorem, CDF characterization, Lebesgue differentiation and elementary monotone-function facts. The source proves the stronger fixed-marginal uniform CDF metric and finite-concatenation continuity; neither is supplied as a free research input.
\item Quantile transform, stochastic order and increasing coordinate homeomorphisms R to(0,1). Interpret the author bounded-coordinate continuity proof for general real marginals. No common-time approximation theorem or atomic distributivity theorem is supplied as background.
\end{itemize}\medskip
\ifdefined\hypertarget\hypertarget{sourcepointer-1}{}\fi
\noindent\textit{Source reference.} Original Theorem A: sources/markov-quantile-jep-2022-177/source0.tex, source lines 295--325. Historical existence/minimality theorem; the exact level-kernel construction required for Theorem B is proved in the retained1451-1705 source. No missing global approximation is supplied by this pointer.\par\medskip
\ifdefined\hypertarget\hypertarget{sourcepointer-2}{}\fi
\noindent\textit{Source reference.} Original section A(a)(iv): sources/markov-quantile-jep-2022-177/source0.tex, source lines 313--317. Individual pairwise limiting description recalled only as context. The retained construction1644-1705 states and proves that description; the new common-time and full-process obligation is proved separately.\par\medskip
\ifdefined\hypertarget\hypertarget{sourcepointer-3}{}\fi
\noindent\textit{Source reference.} Original section 4.21(c): sources/markov-quantile-jep-2022-177/source0.tex, source lines 1777--1786. Unselected later classification of acceptable time sets. The counted Theorem B uses only the included(a),(b), whose entire proof is retained.\par\medskip
\ifdefined\hypertarget\hypertarget{sourcepointer-4}{}\fi
\noindent\textit{Source reference.} Original Proposition 3.11: sources/markov-quantile-jep-2022-177/source0.tex, source lines 1055--1086. The required almost-sure predicate(b) is explicitly printed1062-1063; its implication to order preservation is printed1079. The unused global extension(c) is not assumed.\par\medskip
\ifdefined\hypertarget\hypertarget{sourcepointer-5}{}\fi
\noindent\textit{Source reference.} Original section 3.11(a): sources/markov-quantile-jep-2022-177/source0.tex, source lines 1058--1059. Order-preservation reading is exactly the consequence proved in the retained1079; this pointer introduces no free local theorem.\par\medskip
\ifdefined\hypertarget\hypertarget{sourcepointer-6}{}\fi
\noindent\textit{Source reference.} Original section 3.11(b): sources/markov-quantile-jep-2022-177/source0.tex, source lines 1061--1063. The complete predicate is directly included1062-1063; only its label lies outside the selected plain-text definition.\par\medskip
\ifdefined\hypertarget\hypertarget{sourcepointer-7}{}\fi
\noindent\textit{Source reference.} Original section 4.8: sources/markov-quantile-jep-2022-177/source0.tex, source lines 1545--1563. Author explicitly identifies the remark and figures as qualitative and skippable1543. The actual level-kernel monotonicity proof1521-1533 and1578-1583 is retained.\par\medskip
\ifdefined\hypertarget\hypertarget{sourcepointer-8}{}\fi
\noindent\textit{Source reference.} Original section 3.9: sources/markov-quantile-jep-2022-177/source0.tex, source lines 1034--1050. General stochastic-supremum explanation; the CDF definition/order and directed existence lemma with complete proof are directly supplied934-1031. No extra supremum step is hidden here.\par\medskip
\ifdefined\hypertarget\hypertarget{sourcepointer-9}{}\fi
\noindent\textit{Source reference.} Original section 3.1: sources/markov-quantile-jep-2022-177/source0.tex, source lines 931--1031. Section pointer only; actual order definitions and required supremum proof are included.\par\medskip
\paragraph{Scope of the simultaneous approximation theorem.}
Let $(\mu_t)_{t\in\mathbb R}$ be an arbitrary family of probability measures on
$\mathbb R$. The quantile law is obtained from one uniform variable by the
coordinate quantile maps. The Markov--quantile law is the Markov law whose
pair marginals are constructed below from the level kernels $L_{[s,t]}$.
Convergence in the counted Theorem B means convergence of every finite-dimensional
marginal, as specified in the author's Reminder following that theorem.
Increasing kernels are used in the author's almost-sure conditional-expectation
sense, not as a demand for a probability-valued monotone extension outside the
support of the starting marginal. All finite time tuples are ordered increasingly.
The supporting excerpt of the detailed theorem includes parts (a) and (b) only;
its further classification of admissible time sets in part (c) is not counted.

\begin{notation}\label{notation:marg} For all measurable space $E$, $\m(E)$\ and $\p(E)$\ are the spaces of measures and probability measures on $E$. If $\Tt'\subset\Tt$, $\pr^{\Tt'}$ is the projection $\prod_{\tau\in\Tt}E_\tau\to\prod_{\tau\in\Tt'}E_\tau$; in case $\Tt=\{\tau_1,\ldots,\tau_m\}$ is finite, $\pr^{\tau_1,\ldots,\tau_m}$ means $\pr^{\{\tau_1,\ldots,\tau_m\}}$. When $P\in\p\bigl(\prod_{\tau\in\Tt}\espace_\tau\bigr)$ and $s<t$, $P^s$ stands for $(\pr^{s})_\#P$ and $P^{s,t}$ for $(\pr^{s,t})_\#P$, and $\ma((\mu_\tau)_{\tau\in\Tt})$ denotes $\{P\in\p\bigl(\prod_{\tau\in\Tt}\espace_\tau\bigr):\forall\tau\in\Tt,(\pr^\tau)_\#P=\mu_\tau\}$. When not otherwise specified, what we call the marginals of $P$ are its marginals $P^s$ on a single factor.
\end{notation}

\begin{reminder}\label{reminder:weak_convergence} A sequence $(P_n)_n$ of (probability) measures on some measurable topological space $E$ \emph{converges weakly} to $P$ if, for all bounded continuous function $f$, $\int f\dd P_n\to\int f\dd P$. For $E=\R^\R$ with the weak topology, this convergence amounts to the weak convergence of all finite marginals.
\end{reminder}

Everywhere $E$, $E'$, $E''$ etc.\ are topological spaces (or sometimes Polish spaces) and $\mathcal{B}(E)$, $\mathcal{B}(E')$ and $\mathcal{B}(E'')$ their Borel $\sigma$-algebras.

\begin{defi}
A probability kernel, or kernel $k$ from $E$ to $E'$ is a map\vspace*{-3pt}\enlargethispage{.5\baselineskip}%
\[
k:E\times \mathcal{B}(E')\to [0,1]
\]
such that $k(x,\cdot)$ is a probability measure on $E'$ for every $x$ in $E$ and $k(\cdot,B)$ is a measurable map for every $B\in {\mathcal B}(E')$.
\end{defi}

Probability kernels are usually interpreted as transition matrices, see Remark \ref{rem:matrix}: after one step a particle at $x$ in $E$ arrives at a random position in $E'$, distributed with respect to $k(x,\cdot)$. We often have that interpretation in mind.

\begin{remnotation}\label{notation:k_P}Every transport plan $P\in\p(E\times E')$ can be disintegrated with respect to its first marginal $P^1\eqdefup(\pr^1)_\#P$ and a kernel that we denote by $k_P$, defined from $E$ to $E'$, so that:\vspace*{-3pt}
\[\iint f(x,y)\,\dd P(x,y)=\int\left(\int f(x,y)\,k_P(x,\dd y)\right)\,\dd P^1(x)\]
for every bounded continuous function $f$. Observe that $x\mto k(x,\cdot)$ is $P^1$-almost surely uniquely determined.
\end{remnotation}

\subsubsection{Composition and action of kernels}\label{ss:compos} Kernels $k$ from $E$ to $E'$ and $k'$ from $E'$ to $E''$ can be composed as follows:\vspace*{-3pt}
\[(k\ldot k')(x,A)=\int_{E'} k'(y,A) k(x,\dd y).\]
Similarly, acting on the right, kernels from $E$ to $E'$ \emph{transport}, or \emph{send} positive measures $\theta$ on $E$ on positive measures on $E'$. Acting on the left, they send (adequately integrable) functions $f:E'\to \R$, on functions $E\to \R$:\vspace*{-3pt}
\[(\theta\ldot k)(A)=\int_{E'}k(y,A)\theta(\dd y)\quad \text{and:}\quad (k\ldot f)(x)=\int_{E'} f(y) k(x,\dd y).\]
\noindent Associativity holds, e.g., $(k\ldot k')\ldot k''=k\ldot(k'\ldot k'')$, and $\theta\ldot(k\ldot f)=(\theta\ldot k)\ldot f$, where the action of measures on functions is the obvious one. This is consistent with the following remark.

\begin{rem}\label{rem:matrix}
We recall the usual interpretation of the composition as matrix multiplication. If $E=\{x_i\}_{i=1}^{n}$, $E'=\{y_j\}_{j=1}^{n'}$, $E''=\{z_k\}_{k=1}^{n''}$ are finite, a measure $\theta\in\m(E)$ is a row vector $(\theta(\{x_i\}))_{i=1}^{n}=(\theta_i)_{i=1}^{n}$, a kernel $k$ from $E$ to $E'$ is a matrix $k=((k_{i,j})_{i=1}^{n})_{j=1}^{n'}$, where $(k_{i,j})_{j=1}^{n'}$ is the measure $k(x_i,\,\cdot\,)\in\m(E')$, viewed as a vector, and a function $f$ from $E''$ to $\R$ is a (column) vector $f=(f(z_j))_{j=1}^{n''}$. Then, taking $k'=((k'_{j,k})_{j=1}^{n'})_{k=1}^{n''}$ a kernel from $E'$ to $E''$, $\theta\ldot k$, $k\ldot k'$ and $k'\ldot f$ introduced above have the same sense as products of matrices.
\end{rem}

\begin{notation}\label{not:ident} We denote by $\idk_E$ the identity kernel (that acts trivially) $(x,B)\mto\delta_x(B)=\one_B(x)$.
\end{notation}

\begin{notation}\label{notation_join} With $\mu\in\m(E)$ and $k$ a kernel from $E$ to $E'$ is naturally associated the law $\join(\mu,k)\in\m(\espace\times \espace')$, having $\mu$ as first marginal and the family $(k(x_0,\cdot))_{x_0\in E}$ as laws (on $E'$) conditioned by $x_0\in E$:\vspace*{-3pt}
\[\forall B,B'\in{\mathcal B}(E)\times{\mathcal B}(E'),\quad \join(\mu,k)(B\times B')=\int_Bk(x,B')\dd\mu(x).\]
In particular, $P=\join(P^1,k_P)$.
\end{notation}

\subsubsection{Composition of transport plans}\label{subsubsec:composition} If $P\in\ma(\mu,\mu')$ and $Q\in\ma(\mu',\mu'')$, we can compose them in a similar way as we compose kernels, getting the product:\vspace*{-3pt}
\[P\ldot Q\eqdefup\join(\mu,k_P\ldot k_Q)\in\ma(\mu,\mu''),\quad \text{so that: }k_{P\ldot Q}=k_P\ldot k_Q.\]

\begin{notation}\label{notation:transport_id}
We denote $((\Id_E)_{i=1}^n)_\#\mu\in\ma((\mu)_{i=1}^n)$ by $\Id_{n,\mu}$ or simply $\id_n$ when there is no ambiguity. With $n=2$, $\Id_{2,\mu}=\join(\mu,\idk_E)$. It is moreover the identity transport: $\id_{2,\mu}{}\ldot P=P=P\ldot \id_{2,\mu'}$.
\end{notation}

\subsubsection{Action of transport plans on measures and functions}\label{subsubsec:action_des_transports} If $\mu\in\m(E)$ and $\mu'\in\m(E')$, transport plans $P\in\ma(\mu,\mu')$ have an action similar to that of kernels from $E$ to $E'$, on ($\mu'$-almost surely defined) classes of functions $f$ and on measures $\theta\in\m(E;\mu)$ absolutely continuous with respect to $\mu$. For instance, the latter are transported in $\m(E',\mu')$, as follows:
\[\text{if $\theta\ll \mu$ has density $g$ and $B'\in{\mathcal B}(E')$,}\quad(\theta\ldot P)(B')=\int_{B'}\int_E g(x)\dd P(x,y).\]
If $k$ is a kernel from $E$ to $E'$, $\theta\ldot k=\theta\ldot \join(\mu,k)$. Conversely $\theta\ldot P=\theta\ldot k_P$.

\begin{defi}\label{defi:transpose} Take $P\in\ma(\mu_1,\ldots,\mu_k)$, or $P\in\ma((\mu_t)_{t\in\R})$. We define its transpose $\trans P\in\ma(\mu_k,\ldots,\mu_1)$, \resp $\trans P\in\ma((\mu_{-t})_{t\in\R})$ by: $ \trans P(B_1\times \ldots\times B_k)=P(B_k\times\ldots\times B_1)$, \resp $\trans P(\prod_iB_{t_i})=P(\prod_iB_{-t_i})$.
\end{defi}

The notation $\trans P$ is a reference to a transposed matrix in Remark \ref{rem:matrix}. For $P=\ma(\mu,\nu)$ we will often consider the bilinear map
\[B:(\theta,f)\mto (\theta P)f=\theta (Pf)=\iint g(x)f(y)\,P(\dd x,\dd y),\]
where $g$ is the density of $\theta$ with respect to $\mu$. The case $\theta=\mu\lfloor_{\op]-\infty,x]}$, $f=\one_{\op]-\infty,y]}$ with $B(\theta,\nu)=P(\op]-\infty,x]\times\op]-\infty,y])$ is of special interest, see Section~ \ref{subsec:stochastic_order}.

\subsubsection{Concatenation of transport plans}\label{subsec:caten} (See, e.g., \cite[p.\,111]{Ke72}, \cite[p.\,23]{Vi2}.)
\begin{defi}\label{defi:concatenation}
If $\mu_i\in\p(E_i)$ for $i\in\{1,2,3\}$, if $P_{1,2}\in\ma(\mu_1,\mu_2)$ and $P_{2,3}\in\ma(\mu_2,\mu_3)$, their concatenation $P_{1,2}\circ P_{2,3}$ is the unique $R\in\p(\R^3)$ such that for every $ (B_1,B_2,B_3)\in\mathcal{B}(E_1)\times\mathcal{B}(E_2)\times\mathcal{B}(E_3)$:
\begin{align}\label{eq:concatenation}
R(B_1\times B_2\times B_3)&=\int_{x\in B_1}\int_{y\in B_2}\int_{z\in B_3} \dd\mu_1(x)k_{1,2}(x,\dd y)k_{2,3}(y,\dd z).
\end{align}
In particular, $R\in\ma((\mu_1,\mu_2,\mu_3)$, $(\pr^{1,2})_\#R=P_{1,2}$, and $(\pr^{2,3})_\#R=P_{2,3}$.
\end{defi}

\begin{rem} Let $k_{i,j}$ be a disintegration kernel for $P_{i,j}$ and $P_{2,1}\eqdefup\trans P_{1,2}$. The right side of \eqref{eq:concatenation} also reads: $\iint_{B_1\times B_2} \int_{B_3}k(y,\dd z)\dd P_{1,2}(x,y)$, hence also:
\begin{equation}
\int_{y\in B_2} \iint_{x\in B_1,\,z\in B_3}k_{2,1}(y,\dd x)k_{2,3}(y,\dd z)\mu_2(\dd y).\label{eq:concatenation3}
\end{equation}
Concatenation is ``reversed when time is reversed'' in the sense that $\trans P_{2,3}\circ{}\trans P_{1,2}=\trans(P_{1,2}\circ P_{2,3})$; this is immediate after \eqref{eq:concatenation3}. Formula \eqref{eq:concatenation} gives immediately that $\circ$ is associative, leading to its following generalization:
\[P_{1,2}\circ\ldots\circ P_{n-1,n}\Bigl({\prod_{i=1}^n}B_i\Bigr)=\int_{(x_i)_i\in \prod_iB_i}
\hspace{-4ex}\dd\mu_1(x_1)k_{1,2}(x_1,\dd x_2)\ldots k_{n-1,n}(x_{n-1},\dd x_n).\]
\end{rem}

\skpt
\begin{rem}
\begin{enumerate}
\item
In Section~\ref{subsubsec:composition}, one can also define $P\ldot Q$ as $\pr^{E,E''}_\#(P\circ Q)$.

\item
The composition and concatenation of transport plans find an easy interpretation in terms of random variables. If $(X_1,X_2,X_3)$ is a random vector with $\law(X_1,X_2)=P$ and $\law(X_2,X_3)=Q$ such that $(X_i)_{i\in\{1,2,3\}}$ is a Markov process ($X_1$ and $X_3$ are independent conditionally on $X_2$), then $P\ldot Q$ is the law of $(X_1,X_3)$ and $P\circ Q$ the law of $(X_1,X_2,X_3)$, see \eqref{eq:concatenation3}.
\end{enumerate}
\end{rem}

\subsubsection{The Markov Property}

We introduce the Markov property here in an alternative, equivalent way to the usual one.

\begin{remdefi}\label{defi:markov2}As recalled in \eqref{eq:def_markovien}, a process $(X_t)_{t\in\R}$ is said to be \emph{Markov}~if:\vspace*{-3pt}
\[
\forall s\in\R,\forall t>s,\quad \law(X_t|\,(X_{u})_{u\leqslant s})=\law(X_t|\,X_{s}).
\]
Denoting $\law((X_t)_t)\in\p(\R^\R)$ by $P$, this is equivalent to the fact that for all finite subset $S=\{s_1,\ldots,s_d\}$ of $\R$, $(\pr^S)_\#P$ is the concatenation $P^{s_1,s_2}\circ\cdots\circ P^{s_{d-1},s_d}$, where $P^{s_i,s_{i+1}}$ denotes $(\pr^{\{s_i,s_{i+1}\}})_\#P$. More generally, we say that any measure $P\in\p(\R^\R)$ is Markov if it satisfies this property.

We extend these definition in the obvious way to processes indexed on subsets $R\subset \R$ and measures on $\p(\R^R)$.
\end{remdefi}

We recall the Kolmogorov--Daniell theorem and its usual corollary on Markov measures.

\begin{pro}[Kolmogorov--Daniell theorem]\label{pro:inductive_limit}
Let $E$ be a Polish space and, for each finite subset $S$ of\/ $\R$, $\mu_S$ be a probability measure on $E^S$. If $(\pr^{S'})_\#\mu_S=\mu_{S'}$ for any $S'\subset S$, there exists a unique $P\in \p(E^\R)$ with $(\pr^S)_\#P=\mu_S$ for all $S$.
\end{pro}

One of the most usual applications of Proposition \ref{pro:inductive_limit} is for measures $\mu_S$ of type $\mu_{s_1,s_2}\circ\cdots\circ \mu_{s_{d-1},s_d}$, where $S=\{s_1,\ldots,s_d\}$.

\begin{cor}\label{cor:consist}
Let $(\mu_{s,t})_{s<t}$ be a family of transport plans in $\p(E\times E)$ such~that:\vspace*{-3pt}
\[\mu_{s,u}=\mu_{s,t}{}\ldot\mu_{t,u}\]
for every $s<t<u$. Then there exists a unique Markov measure $P\in \p(E^\R)$ with $P^{s,t}=\mu_{s,t}$ for every $s<t$.
\end{cor}

\begin{defi}\label{defi:consist}It is usual to call \emph{consistent family} every family $(\mu_S)_S$ or $(\mu_{s,t})_{s<t}$ as in Proposition \ref{pro:inductive_limit} and Corollary \ref{cor:consist}.
\end{defi}

\begin{them}[{\cite[Th.\,1]{Ke72}}]\label{them:existence}
Let $(\mu_t)_{t\in \R}$ be a family of probability measures on some Polish space $E$, and for every $s< t$ let $\mathcal{N}_{s,t}\subset\p(E^2)$ be a set of transport plans.
Assume that:
\begin{enumeratei}
\item\label{them:existence1}
for every $s,t$, $\mathcal{N}_{s,t}$ is not empty,
\item\label{them:existence2}
for every $s,t$, $\mathcal{N}_{s,t}\subset\ma(\mu_s,\mu_t)$,
\item\label{them:existence3}
for every $s,t$, $\mathcal{N}_{s,t}$ is closed for the weak topology,
\item\label{them:existence4}
for $r<s<t$ and any $(P,P')\in\mathcal{N}_{r,s}\times \mathcal{N}_{s,t}$, $P\ldot P'\in\mathcal{N}_{r,t}$,
\item\label{them:existence5}
 for every $d$ and $t_1<\cdots<t_d$, if the sequences $(Q^n_{t_i,t_{i+1}})_n\in \mathcal{N}_{t_i,t_{i+1}}$ converge weakly to $Q_{t_i,t_{i+1}}$, then the sequence $(Q^n_{t_1,t_2}\circ\cdots \circ Q^n_{t_{d-1},t_d})_{n}$ tends weakly to $Q_{t_1,t_2}\circ\cdots \circ Q_{t_{d-1},t_d}$.
\end{enumeratei}

Then, there exists a Markov measure $P\in\ma((\mu_t)_t)$ with $(\pr^{s,t})_\#P \in \mathcal{N}_{s,t}$ for every $s<t$.
\end{them}

\begin{notation}\label{notation:intervallesRd}
\begin{enumerate}
\item
 Let us denote $(x_i)_{i=1}^d$ and $(y_i)_{i=1}^d$ in $\R^d$ by $x$ and $y$. We endow $\R^d$ with the natural partial order defined by:
\[x\leqslant y\quad \text{if:}\quad\forall i,x_i\leqslant y_i.\]

We also set $[x,y]\eqdef\{z\in\R^d:x\leqslant z\leqslant y\}=\prod_i[x_i,y_i]$ and similarly $]x,y]$, etc.\ In~particular $\op]-\infty,x]=\op]-\infty,x_1]\times\cdots\times\op]-\infty,x_d]$.

\item
Several times appear statements where some intervals have to be considered closed or open at some of their bounds, either arbitrarily or depending on possible cases. To alleviate the writing, we introduce the symbol ``$\rsemibracket$'' and place it at these bounds.
\end{enumerate}
\end{notation}

\begin{defi}\label{def:cumul}
If $\mu\in\m(\R^d)$, its cumulative distribution function $F_\mu$, that we also denote by $F[\mu]$ to avoid multiple subscripts, is defined, using Notation \ref{notation:intervallesRd}, by:
\[F_\mu:x\in\R^d\mto\mu(\op]-\infty,x]).\]
\end{defi}

\begin{reminder}\label{remind:F}
Recall for instance from \cite[Th.\,3.25]{Kal} that such functions $F$ are characterized by the fact that:
\begin{enumerate}
\item
\label{item:Fonctions_partielles} for the natural partial order of $\R^d$ (see Notation \ref{notation:intervallesRd}), $F$ is increasing and upper semi-continuous in the sense that for all $x\in\R^d$:
\begin{equation}\label{eq:usc}
\forall\varepsilon>0,\,\exists\eta>0,\,\forall y\geqslant x,\quad \|y-x\|_\infty\leqslant\eta\implies F(x)\leqslant F(y)\leqslant F(x)+\varepsilon.
\end{equation}

\item
\label{item:Flimite} $\lim_{\min_i(x_i)\to-\infty} F(x)=0$ and $\lim_{\min_i(x_i)\to+\infty} F(x)=1$,

\item
\label{item:Fdifferences_croisees} for every $h=(h_1,\ldots,h_d)\in[0,+\infty\clo[^d$ and $x
\in \R^d$, the quantity
\[
\sum_{\eps} \sigma(\eps) F(x+\eps h),
\]
which is the measure of the rectangle $]x,x+h]\subset \R^d$, is non-negative. Here $\eps=(\eps_1,\ldots,\eps_d)$ ranges over $\{0,1\}^d$, $\sigma$ is $1$ if $\sum \eps_i$ is even, $-1$ otherwise, and $\eps h$ means $(\eps_1h_1,\ldots,\eps_dh_d)$.
\end{enumerate}
\end{reminder}

\begin{defi}\label{defi:geqlc} If $d\in\N^\ast$ and $m\in\op]0,+\infty\clo[$, following \cite[\S 6.G]{SS}, we define the \emph{lower orthant order} on $\{\mu\in\m(\R^d):\,\mu(\R^d)=m\}$ by: $\mu\leqlc\nu$ if $F_\mu\geqslant F_\nu$.
\end{defi}

\begin{defi}\label{def:lcsup}We call lower orthant supremum of a family $(P_\tau)_{\tau\in\Tt}$ of measures of same mass $m$ on $\R^d$, the smallest upper bound of $\{P_\tau\} _\tau$ for $\leqlc$, if it exists, \ie a~measure $P$ of mass $m$ such that:
\begin{itemize}
\item
for every $\tau$, $P_\tau\leqlc P$,
\item
$P\leqlc Q$ as soon as $P_\tau\leqlc Q$ for every $\tau$.
\end{itemize}
By definition, if it exists it is unique. We denote it by $\lcsup_\tau P_\tau$. Similarly we define $\lcinf_\tau P_\tau$.
\end{defi}

\label{pt:condexp} For every increasing $h:\R\to[0,1]$, $P\ldot h$ is $\mu$-almost surely increasing, \ie more exactly, there is an increasing $\tilde{h}:\R\to [0,1]$ such that, for every bounded continuous function $g$:
\[\int g(x)h(y) \dd P(x,y)= \int g(x)\tilde{h}(x)\dd \mu(x).\]

\begin{notation}\label{notation:mdec}
If $\mu\in\m(\R)$ we denote by $\m(\mu)$ and $\p(\mu)$ the sets of positive measures, respectively probability measures, absolutely continuous with respect to~$\mu$, and $\mdec(\mu)$ and $\pdec(\mu)$, or $\mdec$ and $\pdec$ if there is no ambiguity, their subsets of measures with bounded and decreasing density.
\end{notation}

\begin{defi}\label{defi:quantile2} Take $\mu\in \p(\R)$. The \emph{quantile function}
\[
G_\mu:[0,1]\to \R\cup\{-\infty,+\infty\}
\]
is the increasing left-continuous function such that $(G_\mu)_\#\Lg=\mu$, \ie the generalized inverse of the cumulative distribution function $F_\mu$:
\[G_\mu(q)=\inf\{x\in\R:\,F_\mu (x)\geq q\}.\]
The quantile measure $\como((\mu_\tau)_{\tau\in\Tt})$ is obtained by pushing forward $\Lg$ on $\p(\R^{\Tt})$ by the map $G=(G_{\mu_\tau})_{\tau\in \Tt}:x\in [0,1]\mto (G_{\mu_{\tau}}(x))_{\tau\in\Tt}\in \R^{\Tt}$.
In other words, the functions~$G_{\mu_{\tau}}$ can be seen as random variables defined on $([0,1],\Lg)$ and $\como((\mu_\tau)_{\tau\in\Tt})$ is the law of the process $(G_{\mu_{\tau}})_{\tau\in\Tt}$. One can check from this definition the equality in~\eqref{eq:HF} and \eqref{eq:HF2}.
\end{defi}

\begin{defi}\label{defi:quantile_discretement_markovien} If $M\in\ma((\mu_t)_t)$ and if $R=\{r_1,\ldots,r_m\}\subset\R$, we denote by $M_{[R]}\in\ma((\mu_t)_{\in\R})$ the measure \emph{$M$ made Markov at the points of $R$} defined by the data of its finite marginals $(\pr^{S})_\# M_{[R]}$, for all finite $S$ containing $R$, as follows:
\[(\pr^{S})_\# M_{[R]}= M^{s^0_1,\ldots, s^0_{n_0},r_1}\circ M^{r_1,s^1_{1},\ldots,s^1_{n_1},r_2}\circ\ldots\circ M^{r_m,s^m_{1},\ldots,s^m_{n_m}},\]
where $S=\{s^0_1,\ldots, s^0_{n_0},r_1,s^1_{1},\ldots,s^1_{n_1},r_2,\ldots,r_m,s^m_{1},\ldots,s^m_{n_m}\}$ and where the first or last term disappears if $n_0$ or $n_m$ is null, respectively. These marginals are consistent in the sense of Definition \ref{defi:consist}. So by Proposition \ref{pro:inductive_limit} this defines $M_{[R]}$. We also commit an abuse of language: $M_{[R]}$ is rather the ``law of a process $X$ of law $M$, made Markov at the points of $R$''.
\end{defi}

By Proposition \ref{pro:compo}, for any $R\subset\R$, $(L_{R\cap[s,t]})_{s<t}\in \ma(\la,\la)$ is also a consistent family, thus again Proposition \ref{pro:inductive_limit} enables us to define the following processes on the set of quantile levels.

\begin{defi}\label{defi:lev}
For all $R\subset\R$ we denote by $\lev_R\in\ma((\la_t)_{t\in\R})$ ($\la_t$ denotes~$\la$ at each $t$) the Markov process with 2-marginals $\lev_R^{s,t}=L_{R\cap[s,t]}$. We call $\lev_\R$ the \emph{level process} attached to $(\mu_t)_{t\in \R}$ and denote it by $\lev$.
\end{defi}

\begin{rem}\label{rem:mini_et_Q}In this subsection, using Definition \ref{defi:quantile_discretement_markovien} with $M=\como$ we obtain measures $\como_{[R]}$ linked with $\lev_R$ as follows. After Remark \ref{rem:comp_mini}, $\como_{[R]}^{s,t}=(G_s\otimes G_t)_\#L_{[s,t]\cap R}$, thus, with $G=(\otimes_{t\in \R} G_{t})$, the measures $G_\#\lev_R$ and $\como_{[R]}$ have the same 2-marginals. They are actually equal, as we prove in the ``claim'' at the beginning of the proof of Theorem \ref{them:deuxieme_sec}\eqref{item:them:deuxieme_b}.
\end{rem}

\begin{equation}\label{eq:def_markovien}
\forall s\in\R,\,\forall t>s,\quad \law(X_t|\,(X_{u})_{u\leqslant s})=\law(X_t|\,X_s),
\end{equation}
where $\law(X_t|\,(X_u)_{u\leq s})$ is the law of $X_t$ conditionally to the $\sigma$-algebra generated by the $X_u$ for $u\leq s$. See also Definition \ref{defi:markov2}, where the Markov property is introduced only through notions defined in this article.

\begin{themBJ}\label{them:b}There is an increasing sequence $(R_n)_{n\in\N}$ of finite subsets of $\R$ such that $\como_{[R_n]}\in\ma((\mu_t)_{t\in\R})$ converges weakly to $\mq$.
\end{themBJ}

\begin{lem}[Continuity of $\circ$ when the kernels are increasing]\label{lem:closed}Take $
(\mu_1,\ldots,\mu_n)\in\p(\R)^n$ and for all $i\in\llbracket1,n-1\rrbracket$ a closed set $\mathcal{I}_{i,i+1}\subset\ma(\mu_i,\mu_{i+1})$ of transport plans with increasing kernel. The sets $\mathcal{I}_{i,i+1}$ satisfy Property \eqref{them:existence5} stated for the sets $\mathcal{N}_{t_i,t_{i+1}}$ in Theorem \ref{them:existence}.
\end{lem}

\begin{lem}\label{lem:transitions_croissantes_limite} Take $\mu$ and $\mu'$ in $\p(\R)$. The space of transport plans with increasing kernel in $\ma(\mu,\mu')$ is closed for the weak topology.
\end{lem}

\begin{rem}\label{rem:lcsup}
\begin{enumerate}
\item
\label{p:lcsupa} In Reminder \ref{remind:F}, \eqref{item:Fonctions_partielles} and the first limit of \eqref{item:Flimite} pass to the infimum of functions that are both monotone and upper semi-continuous. If moreover $(P_\tau)_\tau$ has an upper bound $P$, then the second limit of \eqref{item:Flimite} holds. Indeed, for all $\tau$, $F[P_\tau]\geqslant F[P]$, so~that $\inf_\tau F[P_\tau]\geqslant F[P]$, and besides $\lim_{\min_i(x_i)\to+\infty}F[P](x)=1$. Thus, if $(P_\tau)_\tau$ is bounded from above, $\inf_\tau F[P_\tau]$ satisfies \eqref{item:Fonctions_partielles}--\eqref{item:Flimite} of Reminder \ref{remind:F}, so is a cumulative distribution function if and only if it satisfies \eqref{item:Fdifferences_croisees}.
\item
\label{p:lcsupb} If \eqref{item:Fdifferences_croisees} of Reminder \ref{remind:F} is satisfied by the functions $F[P_\tau]$, then $(P_\tau)_\tau$ has a lower orthant supremum, and $F[\lcsup_\tau P_\tau]=\inf_\tau F[P_\tau]$.
\end{enumerate}
\end{rem}

\begin{remnotation}\label{nota:stosup} If $d=1$ the order $\leqlc$ is usually called stochastic order and denoted by $\leqs$; we will then call ``stochastic supremum'' the lower orthant supremum of Definition \ref{def:lcsup} and denote it by $\stosup$.
\end{remnotation}

\skpt
\begin{lem}[Existence criteria for $\lcsup$]
\label{lem:sto_order}
\begin{enumerate}
\item
\label{item:suite_croissante_majoree} If, for $\leqlc$, a sequence $(P_n)_{n\in \N}\!\in\!(\p(\R^d))^\N$ is bounded from above, and increasing, \ie $n\!\leqslant\! m\Rightarrow P_{n}\!\leqlc\! P_{m}$, then $\lcsup_nP_n$ exists and $(P_n)_{n}$ converges weakly~to~it.

\item
\label{item:treillis} If a family $(P_\tau)_{\tau\in\Tt}$ is bounded from above for $\leqlc$ and if for every $\tau,\tau'\in \Tt$ there exists $\sigma\in \Tt$ such that $P_\sigma\geqlc P_\tau$ and $P_\sigma\geqlc P_{\tau'}$, then $\lcsup_\tau P_\tau$ exists and there is an increasing sequence $(P_{\tau_n})_n$
that converges weakly to it.
\end{enumerate}
The results extend in an obvious way to measures of mass $m>0$ in $\m(\R)$.
\end{lem}

\skpt
\begin{proof}
\eqref{item:suite_croissante_majoree}
Set $F\eqdef\inf_nF[P_n]$. After Remark \ref{rem:lcsup}, showing that $F$ satisfies \eqref{item:Fdifferences_croisees} of Reminder \ref{remind:F} ensures the existence of $\lcsup_nP_n$. Consider $M\eqdef\sum_{\eps} \sigma(\eps) F(x+\eps h)$ as in Reminder \ref{remind:F}\eqref{item:Fdifferences_croisees}. Since $F[P_n]$ is decreasing, $F$ is its simple limit, so $M$ is the simple limit of $\sum_{\eps} \sigma(\eps) F[P_n](x+\eps h)$, hence $M\geqslant0$. We are done. The weak convergence is given by the pointwise convergence of the cumulative distribution functions, see Reminder~\ref{remind:billingsley}.

\eqref{item:treillis}
This is a diagonal construction. Let $\mathcal{C}=\bigcup_{k\in\N}\{x_k\}$ be a countable dense set in $\R^d$. Set $F\eqdef\inf_{\tau\in\Tt}F_\tau$. Then for every $(k,n)\in\N^{\ast2}$ we find $\tau_{k,n}\in \Tt$ such that $F[P_{\tau_{k,n}}](x_k)\leqslant F(x_k)+\sfrac1n$. For each $n$, by a finite induction using the assumption of~\eqref{item:treillis} on the $P_{\tau_{k,n}}$, we find $\sigma_{n}\in \Tt$ such that: $\forall k\leqslant n,\, F[P_{\sigma_{n}}](x_k)\leqslant F(x_k)+\sfrac1n$ and $(P_{\sigma_n})_n$ is increasing. Hence:
\begin{equation}\label{eq:convergence_sur_partie_dense}
\forall x\in \mathcal{C},\quad F[P_{\sigma_{n}}](x)\to F(x).
\end{equation}
By \eqref{item:suite_croissante_majoree}, $P=\lcsup_n P_{\sigma_n}$ exists and:
\begin{equation}\label{eq:conv_sur_Rd}
\forall x\in \R^d,\quad F[P_{\sigma_{n}}](x)\to F[P](x).
\end{equation}
Let us prove that \eqref{eq:convergence_sur_partie_dense} holds for any $x\in\R^d$, so that $F=F[P]$. Assume by contradiction that, for some $x$, $F(x)<\inf_nF[P_{\sigma_n}](x)$, \ie by definition of $F$, that for some $\tau\in\Tt$, $F(x)<F[P_\tau](x)<\inf_nF[P_{\sigma_n}](x)=F[P](x)$. Now $F[P_\tau]$ is upper semi-continuous so in a neighbourhood $\mathcal{U}$ of $x$ in $[x,+\infty\clo[$, $F[P_\tau]<F[P]$. But on the dense set $\mathcal{C}$, by~\eqref{eq:convergence_sur_partie_dense} and \eqref{eq:conv_sur_Rd}, $F[P]=F$, hence on $\mathcal{U}\cap\mathcal{C}\neq\varnothing$, $F[P_\tau]<F$, a contradiction.

Finally $F=F[P]$, so $F$ is a cumulative distribution function, so by Remark \ref{rem:lcsup}\eqref{p:lcsupb}, $P=\lcsup_{\tau\in\Tt}P_{\tau}$. Moreover, $P_{\sigma_n}\to P$.
\end{proof}

\eqref{prodef:transitionsb} $\Rightarrow$ \eqref{prodef:transitionsa}: if $\theta\leqs \theta'$, take $h:\R\to[0,1]$ increasing, then $(\theta \ldot P)\ldot h= \theta \ldot\tilde h\leq \theta' \ldot\tilde h$ since $\theta\leqs \theta'$ and $\tilde h$ is increasing by \eqref{prodef:transitionsb}. Then $\theta'\ldot\tilde h=(\theta'\ldot P)\ldot h$ yields \eqref{prodef:transitionsa}.

\begin{rem}\label{rem:transitions_croissantes_composee} A direct consequence of the first point of Proposition \ref{prodef:transitions} is that, if $P\in\ma(\mu_1,\mu_2)$ and $Q\in\ma(\mu_2,\mu_3)$ have increasing kernel, so has the result $P\ldot Q$ of their composition.
\end{rem}

\begin{rem}[$\mdec(\mu)$ is closed for the weak topology]\label{rem:mdec} A measure $\theta$ belongs to $\mdec(\mu)$ if and only if its density is bounded and, for all $(a,b,c,d)\in\R^4$:
\begin{equation}\label{eq:reformulation_mdec}
\bigl(a<b<c<d\ \text{and}\ \mu([a,b])\ldot\mu([c,d])\neq0\bigr)\implies \frac{\theta([a,b])}{\mu([a,b])}\geqslant \frac{\theta([c,d])}{\mu([c,d])}.
\end{equation}
``Only if'' is clear. For the ``if'' part, by the Lebesgue differentiation theorem,
\[
x\mto\lim_{\eps\to0}\frac{\theta([x-\varepsilon,x+\varepsilon])}{\mu([x-\varepsilon,x+\varepsilon])}
\]
provides a representative of the density. Now if a sequence $(\theta_n)_n$ satisfies \eqref{eq:reformulation_mdec} and weakly tends (see Reminder \ref{remind:billingsley}) to $\theta\in\m(\mu)$, $\theta$ satisfies it also (if $a$, $b$, $c$ or $d$ is an atom of $\theta$, re-obtain \eqref{eq:reformulation_mdec} by limit of larger intervals).
\end{rem}

\begin{rem}[$P$ has increasing kernel if and only if $\trans\!P$ maps $\mdec(\nu)$ to $\mdec(\mu)$]\label{rem:stable}
In Proposition \ref{prodef:transitions}, (b) is equivalent to the same statement with decreasing functions; in turn, transposing, this means that, for all decreasing $h:\R\to[0,1]$, $(h\nu)\ldot\trans\!P$, which is equal to $(P\ldot h)\ldot\mu$, has decreasing density.
\end{rem}

\subsection{Quantile measures and minimal couplings}\label{ss:minimality}
We define the quantile coupling (Definition \ref{defi:quantile1}) and quantile process law (Definition \ref{defi:quantile_process}) through a minimality property that is crucial in our paper. We also state the direct and more classical Definition \ref{defi:quantile2} of the quantile measure. Further characterizations of these coupling and process are given throughout the paper and in \cite{BouJui2} where the approach is optimal transport. The reader can refer to \cite{RR, Vi1} for more background. See also the papers \cite{Ju_seminaire, Pass}.

\begin{reminder}Take $\mu$ and $\nu$ in $\m(\R)$. Every $P$ of $\ma(\mu,\nu)$ satisfies the Fréchet--Hoeffding bound:
\begin{equation}\label{eq:HF}
\forall(x,y)\in\R^2,\quad F_P(x,y)\leq \min(F_\mu(x),F_\nu(y)).
\end{equation}
\end{reminder}

\begin{prodefinition}\label{defi:quantile1}There is a unique $P$ such that \eqref{eq:HF} is an equality. We call it the Fréchet--Hoeffding, comonotonic or \emph{quantile coupling} and denote it by $\como(\mu,\nu)$. Said briefly: $\como(\mu,\nu)=\lcinf \ma(\mu,\nu)$.
\end{prodefinition}

The proof follows from Definition \ref{defi:quantile2}, see below.

\skpt
\begin{rem}[Minimality in the language of transport plans]\label{rem:q_is_mini}
\begin{enumerate}
\item
The quantile process $\como$ may also be defined by the fact that its transitions are \emph{minimal} among those of all transport plans of $\ma(\mu,\nu)$: for any $P\in\ma(\mu,\nu)$ and fixed $x$, $F[\mu\lfloor_{\op]-\infty,x]}{}\ldot k_P](y)= F_P(x,y)$ for every $y\in \R$, hence, after Definition~\ref{defi:quantile1} and the characterization of stochastic order of Remark \ref{rem:stosup_existe}:
\begin{equation}\label{eq:minimalite_des_transitions}\forall P\in\ma(\mu,\nu),\quad\mu\lfloor_{\op]-\infty,x]}{}{}\ldot\como(\mu,\nu)\leqs\mu\lfloor_{\op]-\infty,x]}{}{}\ldot k_P
\end{equation}
(in fact, $\mu\lfloor_{\op]-\infty,x]}{}{}\ldot\como(\mu,\nu)$ is of the type $\nu\lfloor_{\op]-\infty,y[}+a\delta_y$), \ie the quantile coupling maps the measures $\mu\lfloor_{\op]-\infty,x]}$ to the stochastically smallest possible measures, \eqref{eq:minimalite_des_transitions}~being also true for any $\theta\in \mdec(\mu)$ in place of $\mu\lfloor_{\op]-\infty,x]}$.

Notice that a property similar to \eqref{eq:minimalite_des_transitions}, with $\leqc$ in place of $\leqs$, defines the (left\nobreakdash-)curtain coupling in \cite{BJ}.

\item
As $\mu\lfloor_{\op]x,+\infty\clo[}{}{}\ldot k_P=\nu-\mu\lfloor_{\op]-\infty,x]}{}{}\ldot k_{P}$, $\mu\lfloor_{\op]x,+\infty\clo[}{}{}\ldot\como(\mu,\nu)$ is paradoxically \emph{maximal} for $\leqs$, hence minimal transitions mean that the mass of $\mu$ is mixed as less as possible when transported on that of $\nu$.
\end{enumerate}
\end{rem}

\begin{prodefinition}\label{defi:quantile_process}If $(\mu_\tau)_{\tau\in\Tt}$ is a family of measures, there is a unique measure $\como\in \ma((\mu_{\tau})_{\tau\in\Tt})$ such that for every $\tau\neq\sigma$, the transport plan $\como^{\tau,\sigma}=(\pr^{\{\tau,\sigma\}})_\#\como$ is the quantile coupling $\como(\mu_\tau,\mu_\sigma)$. We call it the \emph{quantile measure} of $(\mu_\tau)_{\tau\in\Tt}$.
\end{prodefinition}

\begin{proof}[Proof of Propositions \ref{defi:quantile1} and \ref{defi:quantile_process}] The existence parts follows from Definition \ref{defi:quantile2} below and are proved just after it; in Proposition \ref{defi:quantile1} uniqueness is clear; let us prove it in Proposition \ref{defi:quantile_process}. It is rather easy to prove that the equality in \eqref{eq:HF} for every pair of measures implies the equality in the Fréchet--Hoeffding bound of general dimension:
\begin{align}\label{eq:HF2}
\forall(x_i)_{i=1}^d\in\R^d,\quad F_{P^S}(x_1,\ldots,x_d)\leq \min_{i\in\{1\,\ldots,d\}}(F_{\mu_{s_i}}(x_i)),
\end{align}
where $S=\{s_1,\ldots,s_d\}$ is any finite subset of $\R$. (Take $j$ such that $F_{\mu_{s_j}}(x_j)=\min_{1\leq i\leq d}(F_{\mu_{s_i}}(x_i))$. For every $k\neq j$,
\[
F[\pr^{\{s_j,s_k\}}_\#P^S](x_j,x_k)=F[\mu_{s_j}](x_j)
\]
so that finally $P^S(\op]-\infty,x])=F[\mu_{s_j}](x_j)$.) Since equalities \eqref{eq:HF2} for all finite $S\subset\R$ are a compatible set of conditions, this proves with Proposition \ref{pro:inductive_limit} that there exists at most one quantile measure $P$ in $\ma((\mu_\tau)_{\tau\in\Tt})$.
\end{proof}

\begin{notation}\label{not:multi} We denote by $\como(\mu_{s_1},\ldots,\mu_{s_d})$ and $\como((\mu_t)_{t\in \R})$ the multidimensional quantile coupling and the quantile measure.
\end{notation}

\begin{rem}[Products of quantile couplings have increasing kernel and map~$\mdec$ on~$\mdec$]\label{rem:quantile_transitions_croissantes}
We will see in Section~\ref{subsec:little_t} that $\como(\mu,\nu)$ is a composition of two transport plans with increasing kernel. Therefore by Remark \ref{rem:transitions_croissantes_composee} it has increasing kernel. Since $\trans \como(\mu,\nu)=\como(\nu,\mu)$, Remark \ref{rem:stable} ensures that $\como(\nu,\mu)$ maps $\mdec(\mu)$ to $\mdec(\nu)$.

All this also ensure both properties for products of quantile couplings.
\end{rem}

\Subsection{Distances $\rho$ and $\tilde \rho$}\label{subsec:rho}

\Subsubsection{A distance that metricizes $\ma(\mu,\nu)$}

\begin{reminder}\label{remind:billingsley}We remind the reader of the ``Portmanteau theorem'' (see, e.g., \cite[Th.\,2.1]{Bi}): the weak convergence on some metric space $E$ endowed with its Borel $\sigma$-algebra is (equivalently) defined by:
\begin{equation}\label{eq:defi-weakconvergence}P_n\underset{n\to\infty}{\longrightarrow}P\ \text{if }P_n(R)\underset{n\to\infty}{\longrightarrow}P(R)\ \text{for all $R$ such that }P(\partial R)=0.
\end{equation}
In $\R^n$, it is equivalent to consider in \eqref{eq:defi-weakconvergence} only sets $R$ of the form $\prod_{i=1}^n\op]-\infty,x_i\clo]$, see \cite[Exam.\,2.3 p.\,18]{Bi}.
\end{reminder}

\begin{pronotation}\label{pro:rho} Let $\rho$ be the function defined on $\p(\R^d)^2$ by:
\[\rho(P,Q)=\|F_P-F_Q\|_\infty.\]
(See Definition \ref{def:cumul} for $F_P$ and $F_Q$.) This is a distance. If $\mu_1,\ldots,\mu_d$ are probability measures on $\R$, it induces the weak topology on $\ma((\mu_i)_{1\leq i\leq d})$. More precisely, for any $P\in\ma((\mu_i)_{1\leq i\leq d})$ and any sequence $(P_n)_{n\in\N}$ of elements of $\ma((\mu_i)_{1\leq i\leq d})$, the following are equivalent:
\begin{enumeratei}
\item\label{pro:rhoi}
For all $x=(x_i)_{i=1}^d$, $P\bigl(\partial\bigl( \prod_{i=1}^d\op]-\infty,x_i\clo]\bigr)\bigr)=0$ $\Rightarrow$ $\lim_nF_{P_n}(x)=F_{P}(x)$,

\item\label{pro:rhoii}
For all $x=(x_i)_{i=1}^d$, $\lim_nF_{P_n}(x)=F_{P}(x)$,

\bgroup\renewcommand{\theenumi}{\roman{enumi}'}\setcounter{enumi}{1}
\item\label{pro:rhoii'}
For all $x=(x_i)_{i=1}^d$, $\lim_n\bigl(P_n\bigl(\prod_i\op]-\infty,x_i\rsemibracket\bigr)\bigr)=P\bigl(\prod_i\op]-\infty,x_i\rsemibracket\bigr)$,
\egroup

\item\label{pro:rhoiii}
$(F_{P_n})_{n\in \N}$ converges uniformly to $F_P$.
\end{enumeratei}
\end{pronotation}

\begin{proof} The fact that $\rho$ is a distance is immediate. To get \eqref{pro:rhoiii} $\Rightarrow$ \eqref{pro:rhoii'} use that $\op]-\infty,x\clo[=\bigcup_{y<x}\op]-\infty,y]$. Now it suffices to prove \eqref{pro:rhoi} $\Rightarrow$ \eqref{pro:rhoii} $\Rightarrow$ \eqref{pro:rhoiii}.

\eqref{pro:rhoi} $\Rightarrow$ \eqref{pro:rhoii}. Take $Q\in\ma((\mu_i)_{1\leq i\leq d})$ and $b=(b_i)_{i=1}^d$. For $b'=(b'_i)_{i=1}^d\geq b$:
\begin{align}\label{eq:rectangles}
F_Q(b')-F_Q(b)=Q(R')-Q(R)\leq \mu_1(\op]b_1,b'_1])+\cdots+ \mu_d(\op]b_d,b'_d]),
\end{align}
where $R'=\op]-\infty,b']$ and $R=\op]-\infty,b]$, which shows that $(F_Q)_Q$ is ``equicontinuous on the right''. Take $\eps>0$. Since each $b_i\in \R$ can be approached from the right by a sequence of non-atomic points for $\mu_i$, there exists $b'\geqslant b$ such that
\[
\mu_1(\op]b_1,b'_1])+\cdots+ \mu_d(\op]b_d,b'_d])<\eps\AND \mu_1(b'_1)+\cdots+ \mu_d(b'_d)=0,
\]
hence in particular $Q(\partial R')=0$ for every $Q\in \ma((\mu_i)_{1\leq i\leq d})$ so that \eqref{pro:rhoi} applies to~$R'$ for $P$ and all the $P_n$. Thus:\vspace*{-3pt}
\begin{align*}
|P_n(R)- P(R)|&\leq|P_n(R)- P_n(R')|+|P_n(R')-P(R')|+|P(R')- P(R)|\\
&\leq 2\eps+|P_n(R')-P(R')|,
\end{align*}
and if \eqref{pro:rhoi} holds, $|P_n(R')-P(R')|\leqslant\eps$ for $n$ great enough, hence \eqref{pro:rhoii} follows.

\eqref{pro:rhoii} $\Rightarrow$ \eqref{pro:rhoiii}. We apply an adapted version of the prior argument. Suppose \eqref{pro:rhoii} and take $\eps>0$. For every $i$ there exists a finite sequence\vspace*{-3pt}
\[
-\infty=b_0^{(i)}(\eps)<\cdots< b^{(i)}_{N(\eps,i)}(\eps)=+\infty
\]
avoiding the big atoms of $\mu_i$, \ie so that $\mu_i(\op]b_k^{(i)},b_{k+1}^{(i)}\clo[)< \eps$ (this is classical and is proved, e.g., in \cite[\S 12]{Bi}, which deals with the modulus of continuity of càdlàg paths). Every $R=\op]-\infty,b]$ contains some rectangle $R^-_\eps$ and is included in the interior of some rectangle $R^+_\eps$, both bounded by consecutive points $(b^{(i)}_k)_{i,k}$. Using again \eqref{eq:rectangles} for the first and last terms:
\begin{align*}
|P_n(R)- P(R)|&\leq\underbrace{|P_n(R)- P_n(R_\eps^-)|}_{\leqslant d\eps}+\underbrace{|P_n(R_\eps^-)-P(R_\eps^-)|}_{\underset{n\to\infty}\longrightarrow0\ \text{by \eqref{pro:rhoii}}\quad\text{\bf$(\ast)$}}+\underbrace{|P(R_\eps^-)- P(R)|}_{\leqslant d\eps},
\end{align*}
where {\bf$(\ast)$} is uniform as the rectangles $R_\eps^-$ are finitely many. We get \eqref{pro:rhoiii}.
\end{proof}

\subsubsection{Another expression for $\rho$; an alternative distance $\tilde \rho$}\label{ss:another}

Let $\mu_1,\ldots,\mu_d$ be probability measures on $\R$ and $P$, $Q$ stand for elements of $\ma((\mu_i)_{i=1}^d)$.

\begin{notation} When $f$, $g$ and $(f_i)_{i=1}^d$ are functions, by a slight abuse in this subsection, $f\ldot g$ denotes the function $(x,y)\mto f(x)g(y)$ and\/ $\prod_{i=1}^df_i$ the function $(x_1,\ldots ,x_d)\mto \prod_{i=1}^df_i(x_i)$
\end{notation}

\begin{rem}\label{rem:product}By definition, $\rho(P,Q)=\sup\bigl|\int \prod_{i=1}^df_i \,\dd P-\int \prod_{i=1}^df_i \,\dd Q \bigr|$, where each $f_i$ ranges over $\{\one_{\op]-\infty,x]}:\,x\in\R\}$. But in fact:
\end{rem}

\skpt
\begin{prodefinition}\label{pro:product}
\begin{enumerate}
\item\label{pro:producta}
For any $P$ and $Q$ in $\ma((\mu_i)_{i=1}^d)$:\vspace*{-3pt}
\begin{equation}\label{eq:rho_comme_sup}
\rho(P,Q)=\sup\biggl|\int\prod_{i=1}^df_i \,\dd P-\int \prod_{i=1}^df_i \,\dd Q \biggr|,
\end{equation}
where each $f_i$ ranges over $\{f:\R\to[0,1]: f\text{ is decreasing}\}$.

\item\label{pro:productb}
Proposition \ref{pro:rho} may be stated with a distance $\tilde\rho$ based on $x\in \R^n\mto P([x,+\infty\clo[)$ in place of $F[P]$. Then $\tilde\rho$ satisfies \eqref{eq:rho_comme_sup} with \emph{increasing} functions $f_i$.
\end{enumerate}
\end{prodefinition}

\begin{proof}
Once \eqref{pro:producta} is shown, \eqref{pro:productb} is clear. To show \eqref{pro:producta}, by Remark \ref{rem:product}, we have only to prove $\geqslant$. For $f:\R\to [0,1]$ a decreasing function and $I(t)\eqdef f^{-1}([t,+\infty\clo[)$, $f(x)=\int_0^1 \one_{I(t)}(x)\dd t$, thus for $(f_i)_{i=1}^d$ such functions:
\[\biggl(\prod_{i=1}^df_i\biggr) (x_1,\ldots, x_d)=\int_{t\in[0,1]^d}\one_{R(t)}(x_1,\ldots,x_d)\dd t,\]
where $R(t)=I(t_1)\times\cdots\times I(t_d)$. Therefore:
\begin{align*}
\biggl|\int \prod_{i=1}^d f_i\,\dd P-\int \prod_{i=1}^d f_i \,\dd Q\biggr|&=\biggl|\int_{t\in[0,1]^d}P(R(t))-Q(R(t))\dd t\biggr|\\
&\leq \int_{t\in[0,1]^d}|P(R(t))-Q(R(t))|\dd t\\
&\leq \rho(P,Q).\qedhere
\end{align*}
\end{proof}

Proposition \ref{pro:product} has a corollary in the case of transport plans ($d=2$).
\begin{pro}\label{pro:multiplicationcontractante}Take $\mu$, $\mu'$ in $\p(\R)$ and $P$, $Q$ in $\ma(\mu,\mu')$. Then:
\[\rho(P,Q)=\sup\{\rho(\theta\ldot P, \theta \ldot Q):\theta\in\mdec(\mu)\ \text{and $\theta$ has density bounded by $1$}\},\]
where the distance $\rho$ on the right of the equality is that on $\p(\R)$.

Hence if, for some $\nu\in \p(\R)$, $R\in\ma(\nu,\mu)$ is a transport preserving $\mdec$, \ie $\forall\eta\in\mdec(\nu)$, $\eta\ldot R\in\mdec(\mu)$, then $\rho(R\ldot P,R\ldot Q)\leqslant\rho(P,Q)$.
\end{pro}

\begin{proof}
For the first equality, as $\int f\ldot g\, \dd P=((f\mu)\ldot P)\ldot g$, by Proposition \ref{pro:product}: $\rho(P,Q)=\sup_{f,g} |((f\mu)\ldot P)\ldot g-((f\mu)\ldot Q)\ldot g|$ and $\rho(\theta\ldot P, \theta \ldot Q)=\sup_g |(\theta\ldot P)\ldot g-(\theta\ldot Q)\ldot g|$. Now $\theta$ is as in the proposition if and only if $\theta=f\mu$ with $f:\R\to[0,1]$, decreasing. The result follows. Then if $R\in\ma(\nu,\mu)$ is as claimed,
\[
\rho(RP,RQ)=\sup\{\rho(\theta\ldot P, \theta \ldot Q)|\,\theta=\bar\theta\ldot R,
\]
where $\bar\theta\in\mdec(\mu)$ and $\bar\theta$ has density bounded by 1$\}$. This set is included in that of the proposition, since the action of $R$ on $\bar\theta$ does not increase the maximum of its density.
\end{proof}

\begin{rem}[$\{P\in\ma(\mu,\nu):P\text{ maps } \mdec(\mu)\text{ to }\mdec(\nu)\}$ is closed for the weak topology]\label{rem:preserver_mdec_ferme}
Any $\theta\in\mdec(\mu)$ is an increasing limit of positive combinations of characteristic functions $\one_{\op]-\infty,x]}$, so $P\in\ma(\mu,\nu)$ maps $\mdec(\mu)$ in $\mdec(\nu)$ if and only if it maps $\{\mu\lfloor_{\op]-\infty,x]}:\,x\in\R\}$ in it. Now take a sequence $(P_n)_n\in\ma(\mu,\nu)^\N$ of transport plans having this property and converging weakly to $P\in\ma(\mu,\nu)$. After Proposition \ref{pro:rho}, $\|F[P_n]-F[P]\|_\infty\to0$; in particular, for any $x\in\R$, $\|F[P_n](x,\,\cdot\,)-F[P](x,\,\cdot\,)\|_\infty\to0$, \ie $\|F[\mu\lfloor_{\op]-\infty,x]}{}\ldot P_n]-F[\mu\lfloor_{\op]-\infty,x]}{}\ldot P]\|_\infty\to0$, \ie $\mu\lfloor_{\op]-\infty,x]}{}\ldot P_n$ converges weakly to $\mu\lfloor_{\op]-\infty,x]}{}\ldot P$. By Remark \ref{rem:mdec}, $\mu\lfloor_{\op]-\infty,x]}{}\ldot P\in\mdec(\nu)$, we are done.
\end{rem}

Now the little Lemma \ref{lem:transitions_croissantes_limite}, which we use several times, is immediate.

\begin{proof}[Proof of Lemma \ref{lem:transitions_croissantes_limite}]\label{proof:lem:transitions_croissantes_limite} Apply Remarks \ref{rem:preserver_mdec_ferme} and \ref{rem:stable}.
\end{proof}

\begin{rem}\label{rem:diese_contractant} Take $(H_i)_{i=1}^d$ increasing real functions, and $H=H_1\otimes\cdots\otimes H_d:\R^d\to\R^d$. Then $H_\#:\ma((\mu_i)_{i=1}^d )\to\ma((H_i(\mu_i))_{i=1}^d)$ is contracting for $\rho$. To~check it, think that $(F[H_\#P]-F[H_\#Q])(x_1,\ldots,x_d)=(H_\#P\linebreak[2]-H_\#Q)(\prod_i\op]-\infty,x_i])$ equals a term of the type $(H_\#P-H_\#Q)(\prod_i\op]-\infty,y_i\rsemibracket)$, and use Proposition \ref{pro:rho}. Thus if $(P_n)_n$ converges to $P$ in $\ma(\mu_1,\ldots,\mu_d)$, $(H_\#P_n)_n$ converges to $H_\#P$ in $\ma(H_1(\mu_1),\ldots, H_d(\mu_d))$.
\end{rem}

The next remark is neither related to Proposition \ref{pro:rho} nor to Lemma \ref{lem:transitions_croissantes_limite} but is an analogue of Remark \ref{rem:diese_contractant} with $\leqs$ in place of $\rho$.

\begin{rem}\label{rem:leqlc_et_Gs}If $P,Q$ are in $\p(\R^d)$, the $(H_i)_{i=1}^d$ are increasing functions, and $H\eqdef H_1\otimes\ldots\otimes H_d$ then $H_\#:\p(\R^d)\to\p(\R^d)$ is increasing for $\leqlc$, \ie $P\leqlc Q\Rightarrow H_\#P\leqlc H_\#Q$. To check it, think that\vspace*{-3pt}
\[\textstyle
(F[H_\#P]-F[H_\#Q])(x_1,\ldots,x_d)=(H_\#P-H_\#Q)(\prod_i\op]-\infty,x_i])
\]
equals a term of the type $(H_\#P-H_\#Q)(\prod_i\op]-\infty,y_i\rsemibracket)$, and, for the indices $i$ such that ``$\rsemibracket$'' stands for ``$[$'', use that $\op]-\infty,x_i\clo[=\bigcup_{n}\op]-\infty,x_i-1/n]$.
\end{rem}

\begin{rem}\label{rem:multiplicationcontractante} Notice then that $\rho(T,T')=\rho(\trans T,{\trans T'})$ (see Definition \ref{defi:transpose} for $\trans T$). Thus, in Proposition \ref{pro:multiplicationcontractante}, if $R\in\ma(\mu',\nu)$ and if $R$ has increasing kernel (\ie by Remark \ref{rem:stable}, $\trans\!R$ maps $\mdec(\nu)$ into $\mdec(\mu')$):\vspace*{-3pt}
\[\rho(PR,QR)=\rho(\trans(PR),{\trans (QR)})=\rho(\trans\!R\,\trans\!P),{\trans\!R\,\trans\!Q)})\leqslant\rho(\trans\!P,\trans Q)=\rho(P,Q).\]
\end{rem}

\subsubsection{Proof of Lemma \ref{lem:closed}}\label{subsec:MaInProcess}
First we prove Lemma \ref{lem:pointwise}, then its consequence Proposition \ref{pro:continuity}, and finally Lemma \ref{lem:closed}.

\begin{lem}\label{lem:pointwise}
For every $n\in \N$, let $P_n\in \ma(\mu,\nu)$ have increasing kernel. Suppose moreover that $(P_n)_{n\in\N}$ converges to $P_0$. Let $h:[0,1]\to [0,1]$ be an increasing function and for every $n\in \N$, $\tilde{h}_n=P_{n}{}\ldot h$. Then $(\tilde h_n)_{n\in\N}$ converges to $\tilde h_0$, $\mu$-almost surely.
\end{lem}
\begin{proof}By the equivalence proved in Proposition \ref{prodef:transitions}, the sequence $(\tilde h_n)_{n\in \N}$ is increasing. By Propositions \ref{pro:product}\eqref{pro:productb} and \ref{pro:rho}, $\tilde{\rho}(P_n,P_0)\to_{n\to\infty} 0$, \ie for every increasing $g$ with values in $[0,1]$, $\int g(y)h(z) \dd P_n(y,z)\to \int g(y)h(z)\dd P_0(y,z)$. Hence:\vspace*{-3pt}%
\begin{equation}\label{eq:lem:pointwise}
\int g(x)\tilde{h}_n(x)\,\dd \mu(x)\to \int g(x)\tilde{h}_0(x)\,\dd \mu(x),
\end{equation}
by definition of $\tilde h_n$. This also holds if $g$ is the difference of two increasing functions, in particular $g=\one_{[a,b]}=\one_{[a,+\infty\clo[}-\one_{[b,+\infty\clo[}$.

Let $A$ be the set of the elements $x$ of $[0,1]$ such that $\mu[x,1]\mu[0,x]>0$ and for every $n\in \N$, $\mu\text{-}\mathrm{esssup}_{[0,x]}\tilde{h}_n=\mu\text{-}\mathrm{essinf}_{[x,1]}\tilde{h}_n=\tilde{h}_n(x)$. In case $\mu=\la$, since $\tilde h_n$ is increasing, hence has at most a countable number of discontinuity points, $\mu(A)$ is $1$. This also holds in the general case and is given by the increasing functions $\tilde{h}_n\circ G_\mu$, we leave the details to the reader. Now we take any $x\in A$ and prove $\tilde h_n(x)\to\tilde h_0(x)$. It is enough to prove $\limsup_{n}\tilde{h}_n(x)\leq \tilde{h}_0(x)$, since $\liminf_{n}\tilde{h}_n(x)\leq \tilde{h}_0(x)$ can be proved symmetrically. Suppose, for contradiction:\vspace*{-3pt}
\[\limsup_{n}\tilde{h}_n(x)\geq \tilde{h}_0(x)+\eps,\quad \text{for some $\eps>0$.}\]
As $\tilde{h}_{0}$ is increasing and $\tilde{h}_0(x)=\mu\text{-}\mathrm{essinf}_{[x,1]}\tilde{h}_0$ there exists $y>x$ such that $\mu[x,y]>0$ and:\vspace*{-3pt}
\[\tilde{h}_0(x)\leq\frac1{\mu[x,y]} \int_{[x,y]}\tilde{h}_0\,\dd\mu\leq \tilde{h}_0(x)+\eps/2.\]
This contradicts the facts that $\frac1{\mu[x,y]} \int_{[x,y]}\tilde{h}_n\,\dd\mu\to_{n\to\infty}\frac1{\mu[x,y]} \int_{[x,y]}\tilde{h}_0\,\dd\mu$, obtained with $g=\one_{[x,y]}$ in \eqref{eq:lem:pointwise}, and that $\tilde h_n$ is increasing.
\end{proof}

\begin{pro}\label{pro:continuity}
Let $P_n$ tend to $P_0$ in $\ma(\mu_1,\ldots,\mu_d,\eta)$ and $P_n'$ tend to $P_0'$ in $\ma(\eta,\nu)$. Assume moreover that $P'_n$ has increasing kernel for every $n$. Then $P_n'\circ P_n$ tends to $P_0'\circ P_0$.
\end{pro}

\begin{proof}
By Propositions \ref{pro:product}\eqref{pro:productb} and \ref{pro:rho}, we must show that\vspace*{-3pt}
\[
\int f(x)g(y)h(z)\dd (P_n\circ P'_n)(x,y,z)\to \int f(x)g(y)h(z)\dd (P_0\circ P_0')(x,y,z),
\]
where $f((x_i)_{i=1}^d)=\prod_if_i(x_i)$ and $g$, $h$ and the $f_i$ are any increasing functions from $\R$ to $[0,1]$. For all $n\geqslant0$:
\begin{align*}
\int f(x)g(y)h(z)\dd (P_n\circ P'_n)(x,y,z)&=\int f(x)g(y)\int h(z) k_{P'_n}(y,\dd z)\dd P_n(x,y)\\
&=\int f(x)(g\tilde h_n)(y)\dd P_n,
\end{align*}
where $\tilde h_n=P_{n}{}\ldot h$, by Definition \ref{defi:concatenation} and Section~\ref{subsubsec:action_des_transports}.

By Lemma \ref{lem:pointwise}, $(g\tilde{h}_n)_n$ pointwise converges to $g\tilde{h}_0$, almost surely. Hence by the dominated convergence theorem:
\[\left|\int f(x)(g\tilde h_n)(y)\dd P_0-\int f(x)(g\tilde h_0)(y)\dd P_0\right|\to0.\]
Moreover, $\tilde{\rho}(P_n,P_0)$ tends to zero and for every $n$, $g\tilde h_n$ is an increasing functions taking values in $[0,1]$. Hence, again by Propositions \ref{pro:product}\eqref{pro:productb} and \ref{pro:rho}:
\begin{equation}\label{eq:triangle}
\left|\int f(x)(g \tilde h_n)(y)\dd P_n-\int f(x)(g \tilde h_n)(y)\dd P_0\right|\to0.
\end{equation}
Therefore with the triangle inequality, \eqref{eq:triangle} holds with $\tilde h_0$ in place of $\tilde h_n$ on the right,
which is exactly what we claimed.
\end{proof}

We can now prove Lemma \ref{lem:closed}. It follows directly from Proposition \ref{pro:continuity}.

\begin{proof}[Proof of Lemma \ref{lem:closed}]\label{proof:lem:closed}
Consider a sequence $P^n_{1,2}\circ\cdots\circ P^n_{d,d+1}$ with limit $P\in\ma(\mu_1,\ldots,\mu_{d+1})$ and denote $\pr^{i,i+1}_{\#}P$ by $P_{i,i+i}$. Since $(\pr^{i,i+1})_\#$ is continuous, $P_{i,i+1}$ is the limit of $(P_{i,i+1}^n)_n$, hence $P_{i,i+1}\in\mathcal{I}_{i,i+1}$. Proposition \ref{pro:continuity} used by induction shows that $P=P_{1,2}\circ\cdots\circ P_{n,n+1}$.
\end{proof}

\subsection{Transitions kernels to and from the space of quantiles \mb$[0,1]$\mn}\label{subsec:little_t}
In this paper we will need to consider the quantile couplings of the measures $\mu_t$ with the reference measure $\la=\Lg$.

\begin{notation}[$q_r$, $k_r$, $\trans k_r$] \label{notation:kernels} For all $r\in R$ we set $q_r= \como(\la,\mu_r)$; thus $\trans q_r$ (see Definition \ref{defi:transpose}) is $\como(\mu_r, \lambda)$. Those couplings admit the respective disintegration kernels:\vspace*{-3pt}
\begin{align*}
k_r:(\alpha,B)\mto \delta_{G_{r}(\alpha)}(B),\ \trans{k}_r(x,\espcdot)=
\begin{cases}
\delta_{F_r(x)}&\!\!\text{if }\mu_r(x)=0\\
(\mu_r(x))^{-1}\la\lfloor_{]F_r(x^-),F_r(x)[}&\!\!\text{if }\mu_r(x)>0.\\
\end{cases}
\end{align*}
\end{notation}

\skpt
\begin{rem}\label{rem:finesse}
\begin{enumerate}
\item
\label{p1:rem:finesse} Notice that $\trans{q}_s \ldot q_t=\join(\mu_s,\trans{k}_s \ldot k_t)=\como^{s,t}$, where $\como^{s,t}=\como(\mu_s,\mu_t)$, see Definition \ref{defi:quantile_process}. Indeed, $\trans{q}_s \ldot q_t=\trans{q}_s \ldot\id_{\la,2}{}\ldot q_t$, then apply \eqref{p2:rem:finesse} below,
which will also be useful farther.

\item
\label{p2:rem:finesse} If some ordered pair $(U,V)$ of variables has law $T\in \ma(\la,\la)$, then\vspace*{-3pt}
\[
\ma(\mu_s,\mu_t)\ni\trans{q}_s \ldot T\ldot q_t=\law(G_s(U),G_t(V))=(G_s\otimes G_t)_\#T.
\]
Indeed, $\trans{q}_s \ldot T\ldot q_t=\pr^{1,4}_\#(\trans{q}_s \circ T\circ q_t)$ and the 4-times process $(G_{\mu_s}(U),U,V,G_{\mu_t}(V))$ has law $\trans{q}_s \circ T\circ q_t$ and is Markov, since the $\sigma$-fields spanned by $U$ and $\{U,G_{\mu_s}(U)\}$ are the same.

\item
\label{p3:rem:finesse} From \eqref{p1:rem:finesse} we get $\trans q_r\ldot q_r=\id_{\mu_r,2}$, so that $\trans q_r\ldot q_r\ldot\trans q_r=\trans q_r$ and $q_r\ldot\trans q_r\ldot q_r=q_r$. \emph{However}, $q_r\ldot\trans q_r\neq \id_{\la,2}$. Indeed, $k_r\ldot\trans{k}_r$ maps any quantile level $\alpha\in \op]0,1\clo[$ on itself \emph{except when} $G_r(\alpha)$ is an atom of $\mu_r$. Actually, $\mu_r$-almost surely:\vspace*{-3pt}
\begin{align*}
k_r\ldot\trans{k}_r(\alpha)=
\begin{cases}
\delta_\alpha&\text{if }\mu_r(G_{r}(\alpha))=0\\
(\alpha^+-\alpha^-)^{-1}\la|_{]\alpha^-,\alpha^+[}&\text{if }\alpha\in]\alpha^-,\alpha^+[,
\end{cases}
\end{align*}
where $]\alpha^-,\alpha^+[$ denotes any set $A_{r,x}$ as follows.
\end{enumerate}
\end{rem}

\skpt
\begin{notation}[Atomic levels, $\ell_r$]\label{nota:Asx}
\begin{enumerate}
\item
\label{p:Asx} We denote by $A_{r,x}$ the interval $\op]F_r(x^-),F_r(x)\clo[$ of quantile levels merged by $G_r$ on some atom $x$ of $\mu_r$, and by $A_r$ the set $\bigcup_{\mu_r(x)>0}A_{r,x}\!\subset\! [0,1]$ of ``atomic levels'' of~$\mu_r$.

\item
\label{p:l_R_fini} We denote $k_r\ldot\trans{k}_r$ described in Remark \ref{rem:finesse}\eqref{p3:rem:finesse} by $\ell_r$.
\end{enumerate}
\end{notation}

\begin{rem}\label{rem:action_l}Measures $\theta\ll \la$ are transported by $\ell_r$ as follows: $\theta\ldot \ell_r$ coincides with~$\theta$ on $\op]0,1\clo[\setminus A_r$, and on each $A_{r,x}$ it has constant density and mass $\theta(A_{r,x})$, \ie equals $(\alpha^+-\alpha^-)^{-1}\theta(A_{r,x})$. Equivalently, $F[\theta\ldot \ell_r]$ is continuous, equal to $F[\theta]$ on $\op]0,1\clo[\setminus A_{r,x}$, and affine on each connected component $A_{r,x}$ of $A_r$.
\end{rem}

The product $\como^{s,r_1}{}\ldot\como^{r_1,r_2}\lcdots \como^{r_{m-1},r_m}{}\ldot\como^{r_m,t}$ appearing in Theorem \ref{them:a} is more deeply analyzed in Theorem \ref{them:b}. Its kernel reads:
\begin{equation}\label{eq:composee_l}
\begin{aligned}
(\trans k_s\ldot k_{r_1})\ldot(\trans k_{r_1}{}\ldot k_{r_2})\lcdots (\trans k_{r_{m-1}}&\ldot k_{r_m})\ldot(\trans k_{r_m}{}\ldot k_t)\\
&= \trans k_s\ldot(k_{r_1}{}\ldot\trans k_{r_1})\lcdots (k_{r_m}{}\ldot\trans k_{r_m})\ldot k_t\\
&=\trans k_s\ldot\ell_{r_1}\lcdots \ell_{r_m}{}\ldot k_t.
\end{aligned}
\end{equation}
In Remark \ref{rem:comp_mini}, \eqref{eq:composee_l} is further commented and re-expressed for transports in place of kernels. For now, it leads to introduce the following kernel.

\begin{notation}\label{notation:Ts}
Take $R=\{r_1,\ldots,r_m\}\!\subset\!\R$. We denote the kernel \hbox{$\ell_{r_1}{}\ldot\ell_{r_2}{}\ldot\cdots\ldot\ell_{r_m}$} from $\op]0,1\clo[$ to itself by $\ell_R$, and $\join(\la;\ell_R)\in\ma(\la,\la)$ by $L_R$. If $R=\varnothing$, $\ell_\varnothing$ is the identity kernel and $L_\varnothing$ the identity transport.
\end{notation}

Notice that $\ell_{\{r\}}=\ell_r$ and that for any $R$, $\la\ldot\ell_R=\la$. Moreover, $\ell_R$ only depends on $(A_r)_{r\in R}$. The following lemma is particularly simple.

\begin{lem}\label{lem:mono1}
Let $R\subset \R$ be a finite set and $\mu,\,\nu$ be in $\m(\la)$.

\begin{enumerate}
\item
\label{item:mono1a} If $\mu\leqs \nu$ then $\mu \ldot\ell_R\leqs \nu\ldot\ell_R$, \ie $\join(\la; \ell_R)$ has increasing kernel.
\item
\label{item:mono1b} If $\mu$ is in $\mdec(\la)$, then so is $\mu \ldot\ell_R$, and $\mu\leqs \mu\ldot\ell_R$.
\end{enumerate}
\end{lem}

\begin{proof} For $r\in \R$, one easily checks that $\ell_r$ is increasing and stabilizes $\mdec$; \eqref{item:mono1a}~and the first point of \eqref{item:mono1b} follow. If $\mu\in \mdec(\la)$, to see that $\mu\leqs \mu\ldot\ell_R$, look at the cumulative distribution functions. They coincide off the components of $A_r$. Now on each of those, $F_{\mu\ldot\ell_R}$ is affine whereas $F_\mu$ is concave since $\mu$ has decreasing density, so necessarily $F_\mu\geqslant F_{\mu\ldot\ell_R}$.
\end{proof}

Though simple, the next lemma is a key of our construction of $\mq$.
\begin{lem}\label{lem:mono2}
Let $R\subset R'$ be two finite subsets of\/ $\R$ and $\mu\in\mdec(\la)$. Then $\mu \ldot\ell_R\leqs \mu \ldot\ell_{R'}$.
\end{lem}
\begin{proof}
Using an induction on the cardinal difference, it is enough to prove this if $R'$ has one more element than $R$, say $r'$. We order it with the elements $r_i$ of $R$: $r_1<\cdots<r_k<r'<r_{k+1}<\cdots<r_m$. By Lemma \ref{lem:mono1}\eqref{item:mono1b}, if $\mu$ is in $\mdec(\la)$, so is $\mu_k\eqdef\mu\ldot\ell_{r_1}\lcdots \ell_{r_k}$ and $\mu_k\leqs \mu_k \ldot\ell_{r'}$. We apply $\ell_{r_{k+1}}\lcdots \ell_{r_m}$ to each term of this inequality. Lemma \ref{lem:mono1}\eqref{item:mono1a} concludes.
\end{proof}

\noindent The choice of the suitable set $\mdec$ in our monotonicity Lemma \ref{lem:mono2} (see also Remark~\ref{rem:qualitative}) enables us to extend the definition of $L_R$ to infinite sets $R$. Recall that $\leqlc$ is interpreted in terms of $\leqs$ in Remark \ref{rem:q_is_mini}.

\skpt
\begin{lemnotation}\label{lem:existenceT}
\begin{enumerate}
\item
\label{p1:lem:existenceT} For all finite subsets $R$, $R^{\prime}$ of\/ $\R$, $L_R\leqlc \la\otimes\la$, and $R'\subset R$ $\Rightarrow$ $L_{R'}\leqlc L_{R}$.

\item
\label{p2:def_L_R} For any $R\subset\R$, $\lcsup\{L_{R'}:\,R'\subset R$ and $R'$ finite$\}$ exists. We denote it by $L_R$ (which is consistent with Notation \ref{notation:Ts} when $R$ is finite). 

\item
\label{p3:def_L_R} For any $R\subset\R$, there is a nested sequence $(R_n)_n$ of finite subsets of $R$ such that $(L_{R_n})_n$ converges weakly to $L_R$; if $(R_n)_n$ satisfies this property, all sequence $(R'_n)_n$ such that $R'_n\supset R_n$ also does.
\end{enumerate}
\end{lemnotation}

\begin{proof} Let us prove \eqref{p1:lem:existenceT}. For all $(x,y)\in[0,1]^2$, $F[L_R](x,y)=(\la\lfloor_{[0,x]})\ldot L_R([0,y])$. After Lemma \ref{lem:mono1}\eqref{item:mono1b}, $(\la\lfloor_{[0,x]})\ldot L_R\in\mdec$ so $f:y\mto\frac1y((\la\lfloor_{[0,x]})\ldot L_R)([0,y])$ is decreasing, thus $f(y)\geq f(1)=x$. Hence:
\[
\forall y\in[0,1],\quad ((\la\lfloor_{[0,x]})\ldot L_R)([0,y])\geqslant xy=F[\la\otimes\la](x,y),\quad \text{\ie} L_R\leqlc \la\otimes\la.
\]
Now if $R'\subset R$ for every $x$ we can apply Lemma \ref{lem:mono2} to $\la\lfloor_{[0,x]}\in\mdec$. This yields $F[L_{R}](x,y)=(\la\lfloor_{[0,x]})\ldot L_{R}([0,y])\geqslant(\la\lfloor_{[0,x]})\ldot L_{R'}([0,y])=F[L_{R'}](x,y)$, which is the expected relation for $\leqlc$. Then, \eqref{p2:def_L_R} and \eqref{p3:def_L_R} follow from criterion \eqref{item:treillis} of Lemma \ref{lem:sto_order}. Indeed, by \eqref{p1:lem:existenceT}, $\mathcal{S}=\{L_{R'}:\,R'\subset R$ and $R'$ finite$\}$ is bounded from above by $\la\otimes\la$, and if $\{L_{R'_1},L_{R'_2}\}\subset\mathcal{S}$, then $L_{R'_1\cup R'_2}\in\mathcal{S}$ and $L_{R'_1\cup R'_2}\geqlc L_{R'_i}$ for $i=1,2$. Finally, the assertion about $(R'_n)_n$ follows from (a) and the interpretation of both $\rho$ and $\leqs$ with cumulative distribution functions.\end{proof}

\begin{notation}\label{nota:kernel_tS}For all $R\subset \R$, we denote by $\ell_R$ the kernel associated with $L_R$. This is again consistent with Notation \ref{nota:Asx}\eqref{p:l_R_fini} when $R$ is finite.
\end{notation}

The following result is crucial to define processes on $(\op]0,1\clo[,\la)$ with Corollary \ref{cor:consist}, as is done in particular in Definition \ref{defi:lev} for $\lev_R$ and $\lev$.
\begin{pro}\label{pro:compo}If $R$ and $R'$ are subsets of $\R$ such that $r\leqslant r'$ for all $(r,r')\in R\times R'$, then
: $L_{R\cup R'}=L_{R}{}\ldot L_{R'}$.
In particular, for $s<t<u$, $L_{[s,u]}=L_{[s,t]}{}\ldot L_{[t,u]}$.
\end{pro}
\begin{proof} Using Lemma \ref{lem:existenceT}\eqref{p2:def_L_R} and \eqref{p3:def_L_R} we find sequences $(R_n)_n$ and $(R'_n)_n$ of finite subsets of $R$ and $R'$ respectively, such that $L_{R_n}$ converges weakly to~$L_R$, $L_{R'_n}$ to~$L_{R'}$, and $L_{R_n\cup R'_n}$ to $L_{R\cup R'}$. Besides, since $r\leqslant r'$ for all $(r,r')\in R\times R'$, since the $R_n$ are finite and since $L_{\{r\}}$ is idempotent (so that in case $R\cap R'\neq\varnothing$ and $R_n\cap R'_n=R\cap R'=\{r\}$, a repetition of $L_{\{r\}}$ does not matter), $L_{R_n\cup R'_n}=L_{R_n}{}\ldot L_{R'_n}$. Then, using Proposition \ref{pro:rho} and the distance $\rho$ introduced in it:
\begin{align*}
\rho(L_R\ldot L_{R'},L_{R\cup R'})&\leqslant\rho(L_R\ldot L_{R'},L_R\ldot L_{R'_n})\\
&\hspace*{2.3cm}+\rho(L_R\ldot L_{R'_n},L_{R_n}{}\ldot L_{R'_n})+\rho(L_{R_n\cup R'_n},L_{R\cup R'})\\
&\leqslant\rho(L_{R'},L_{R'_n})+\rho(L_R,L_{R_n})+\rho(L_{R_n\cup R'_n},L_{R\cup R'})
\end{align*}
by Proposition \ref{pro:multiplicationcontractante} and Remark \ref{rem:multiplicationcontractante}, since $\trans{L_{R'_n}}$ also preserves $\mdec(\la)$. All terms tend to zero when $n$ tends to infinity. The desired equality follows.
\end{proof}

\Subsection{The Markov-quantile process $\mq$; proof of Theorem \ref{them:a}}\label{subsec:preuve_premier}

\begin{reminder} For $(s,t)\in\R^2$ with $t>s$, and $R=\{r_1,\ldots,r_m\}\subset[s,t]$, $\como^{s,t}_{[R]}$ denotes the coupling $\como(\mu_{s},\mu_{r_1}).\como(\mu_{r_2},\mu_{r_3}).\cdots.\como(\mu_{r_{m}},\mu_{t})\in\ma(\mu_s,\mu_t)$, see Theorem \ref{them:a}\eqref{item:limite_de_compose}.
\end{reminder}

\skpt
\begin{prodefinition}\label{pro:comp_conv}
\begin{enumerate}
\item
\label{aaa} The set $\{\como^{s,t}_{[R]}\,|\,R\subset \R\ \text{and}\ \sharp R<\infty\}$ has a lower orthant supremum and there is a nested sequence $(R_n)_{n\in \N}$ of finite sets such that $\como^{s,t}_{[R_n]}$ tends to it. We denote it by $\mu_{s,t}$.

\item
\label{bbb} The family $(\mu_{s,t})_{s<t}$ is consistent in the sense of Definition \ref{defi:consist}, giving rise to a Markov measure $\mq\in\ma((\mu_t)_{t\in\R})$, that we call the \emph{Markov-quantile measure} attached to $(\mu_t)_{t\in\R}$.

\item
\label{ccc} For all $s$ and $t>s$, $\mq^{s,t}=(G_s\otimes G_t)_\#L_{]s,t[}=(G_s\otimes G_t)_\#L_{[s,t]}$.
\end{enumerate}
\end{prodefinition}

To show Proposition \ref{pro:comp_conv} we first state the following crucial relation between $\como^{s,t}_{[R]}\in \ma(\mu_s,\mu_t)$ and $L_R\in \ma(\la,\la)$ defined in Section~\ref{subsec:little_t}.

\begin{rem}\label{rem:comp_mini} Take any finite subset $R$ of $\R$ and $(s,t)$ with $s<t$, then:
\[\como^{s,t}_{[R]}=(G_s\otimes G_t)_\#L_{]s,t[\cap R}=(G_s\otimes G_t)_\#L_{[s,t]\cap R}.\]
Indeed, $\como^{s,t}\!=\!\trans q_s\ldot q_t\!=\!\law(G_s,G_t)$, see Remark \ref{rem:finesse}\eqref{p1:rem:finesse}. It also equals $(G_s\otimes G_t)_\#(\Id_2)$, where $\Id_2=\id_{\la,2}$ is the identity transport form $\la$ to itself, see Notation \ref{notation:transport_id}, since by Remark \ref{rem:finesse}\eqref{p2:rem:finesse}, $\como(\mu_s,\mu_t)=\trans q_s \ldot q_t=\trans q_s. \id_2 \ldot q_t=(G_s\otimes G_t)_\#(\Id_2)$. More generally, for~any $\{r_1,\ldots,r_m\}\subset\op]s,t\clo[$:
\begin{align*}
\como(\mu_{s},\mu_{r_1})\lcdots{}& \como(\mu_{r_m},\mu_{t})\\
&=\trans q_{s}{}\ldot q_{r_1}{}\ldot\trans q_{r_1}{}\ldot q_{r_2}\lcdots \trans q_{r_{m-1}}{}\ldot q_{r_{m}}{}\ldot\trans q_{r_{m}}{}\ldot q_{t}\\
&=(G_{s}\otimes G_{t})_\#(q_{r_1}{}\ldot\trans q_{r_1}{}\ldot q_{r_2}\lcdots \trans q_{r_{m-1}}{}\ldot q_{r_{m}}{}\ldot\trans q_{r_{m}})\\
&\hspace*{2cm}\text{(notice that this writing involves neither $\trans q_{s}$ nor $q_{t}$)}\\
&=(G_{s}\otimes G_t)_\#L_{\{r_1,\ldots,r_m\}}.
\end{align*}

Besides recall that $q_s\ldot\trans q_s\ldot q_s=q_s$ and $\trans q_t\ldot q_t\ldot\trans q_t=\trans q_t$, so that:
\begin{align*}
\como(\mu_{s},\mu_{r_1})\lcdots \como(\mu_{r_m},\mu_{t})
&=(G_{s}\otimes G_{t})_\#(q_{s}{}\ldot\trans q_{s}{}\ldot q_{r_1}\lcdots \trans q_{r_m}{}\ldot q_{t}{}\ldot\trans q_{t})\\
&=(G_{s}\otimes G_t)_\#L_{\{s,r_1,\ldots,r_m,t\}}.
\end{align*}
\end{rem}

\begin{proof}[Proof of Proposition \ref{pro:comp_conv}] We have only to gather our results. We set
\[
\mathcal{S}\eqdef\{R\subset \op]s,t\clo[:\,R\text{ is finite}\}.
\]
By Remark \ref{rem:leqlc_et_Gs} and Lemma \ref{lem:existenceT}\eqref{p1:lem:existenceT},
\begin{align*}
\forall R\in\mathcal{S},\quad (G_s\otimes G_t)_\#L_R&\leqlc (G_s\otimes G_t)_\#(\la\otimes \la),\\
\tag*{and}
R'\subset R\implies(G_s\otimes G_t)_\#L_{R'}&\leqlc (G_s\otimes G_t)_\#L_{R}. 
\end{align*}
Thus, by Lemma \ref{lem:sto_order}\eqref{item:treillis}, $\lcsup_{R'\in\mathcal{S}}(G_s\otimes G_t)_\#L_{R'}$ exists and is the limit of some sequence $((G_s\otimes G_t)_\#L_{R_n})_n$. Hence the limit in \eqref{p1:lem:existenceT} is given by Remark \ref{rem:comp_mini}:
\[\como^{s,r_1}{}\ldot\como^{r_1,r_2}\lcdots \como^{r_{m(n)-1},r_{m(n)}}{}\ldot\como^{r_{m(n)},t}=(G_s\otimes G_t)_\#L_{R_n}.\]

We prove now the first part of (c), \ie $\mu_{s,t}=(G_s\otimes G_t)_\#L_{\op]s,t\clo[}$ and, at the same time, that the sequence $(R_n)_n$ in (a) can be chosen to be nested. Let $(R_n)_n$ be a nested sequence of $\mathcal{S}$ such that $L_{R_n}\to L_{]s,t[}$, \ie $F[L_{R_n}]$ pointwise converges to $F[L_{]s,t[}]$. If $M\in \ma(\mu_s,\mu_t)$ satisfies $(G_s\otimes G_t)_\#L_{R'}\leqlc M$ for every $T\in \mathcal{S}$, this also holds for $R'=R_n$ for all $n$. Now by Remark \ref{rem:diese_contractant}, $(G_s\otimes G_t)_\#L_{R_n}\to (G_s\otimes G_t)_\#L_{]s,t[}$. Therefore, going to the limits at the level of the cumulative distribution functions we~get $(G_s\otimes G_t)_\#L_{]s,t[}\leqlc M$. Then Remark \ref{rem:leqlc_et_Gs} gives
\[
(G_s\otimes G_t)_\#L_{R'}\leqlc(G_s\otimes G_t)_\#L_{]s,t[},
\]
hence:
\[(G_s\otimes G_t)_\#L_{]s,t[}=\lcsup_{R'\subset\R\ \text{and $R'$ finite}}\bigl\{\como^{s,t}_{[R']}=(G_s\otimes G_t)_\#L_{R'}\bigr\}.\]

Corollary \ref{cor:consist} gives \eqref{bbb}. Indeed, Proposition \ref{pro:compo} on the composition of transports~$L_R$ gives the consistency of $(\mu_{s,t})_{s,t}$ (see Definition \ref{defi:consist}):
\[
(G_s\otimes G_t)_\#L_{\op]s,t\clo[}=\trans{q}_s \ldot L_{\op]s,t\clo[}{}\ldot q_t\AND (G_t\otimes G_u)_\#L_{]t,u[}=\trans{q}_t L_{]t,u[}q_u.
\]
Since $q_t\ldot\trans{q}_t=L_{\{t\}}$, $\mu_{s,u}=\mu_{s,t}{}\ldot\mu_{t,u}$.

For the second equality of \eqref{ccc}, proceed as at the end of Remark \ref{rem:comp_mini}.
\end{proof}

\begin{them}\label{them:deuxieme_sec} Let $(\mu_t)_{t\in \R}$ be a family of probability measures on $\R$. In points \eqref{item:them:deuxieme_b} and \eqref{item:them:deuxieme_c}, $(R_n)_n$ stands for a nested sequence of finite subsets of\/ $\R$ and $R$ for $\bigcup_nR_n$.

\begin{enumerate}
\item
\label{item:them:deuxieme_a} There is a countable set $R\subset\R$ satisfying:
\begin{gather}\label{eq:convergence_doug}
\text{for all $(s,t)$ with $s<t$,}\ L_{R\cap\op]s,t\clo[}=L_{\op]s,t\clo[}.
\end{gather}

\item
\label{item:them:deuxieme_b} If $R$ satisfies \eqref{eq:convergence_doug} then (see Reminder \ref{reminder:weak_convergence} for the weak convergence):
\begin{equation}\label{eq:cv}
\text{$(\como_{[R_n]})_n$ converges weakly to $\mq$.}
\end{equation}

\end{enumerate}
\end{them}

\begin{lem}\label{lem:croissance_lc} Let $\Tt$ denote a totally ordered set of indices (in practice, $\Tt=\R^+$ or $\Tt=\N$). If $(R_\tau)_{\tau\in \Tt}$ is a family of subsets of\/ $\R$, increasing for the inclusion, setting $R=\bigcup_\tau R_\tau$, $(L_{R_\tau})_{\tau\in\Tt}$ is increasing for $\leqlc$ and tends weakly to $L_{R}$ when $\tau$ tends to infinity.
\end{lem}

\begin{proof} It rests on this remark: by definition of $\leqlc$ in Definition \ref{defi:geqlc} and of $\rho$ in Proposition \ref{pro:rho}, if $A$, $A'$, $A''$ are measures of $\p(\R^d)$ with the same marginals and $A\leqlc A'\leqlc A''$, then $\rho(A',A'')\leqslant\rho(A,A'')$. Now we prove the lemma. By definition of $L_{R_\tau}$ in Lemma \ref{lem:existenceT}\eqref{p2:def_L_R} the sequence $(L_{R_\tau})_{\tau\in\Tt}$ is increasing, be the sets $R_\tau$ finite or not. By Lemma \ref{lem:existenceT}(c), and as $\rho$ metrizes the weak topology (Proposition \ref{pro:rho}), for all $\eps>0$ we find a finite $R'\subset R$ such that $\rho(L_R,L_{R'})\leqslant\eps$. As $R=\bigcup_\tau R_\tau$ there is a $\tau_0$ such that $R_{\tau_0}\supset R'$. Then: $\tau\geqslant \tau_0\Rightarrow L_{R'}\leqlc L_{R_\tau}\leqlc L_{R}\Rightarrow \rho(L_{R'},L_{R_\tau})\leqslant\eps$ by the remark.
\end{proof}

\begin{lem}\label{lem:convergence_couples}
There is a nested sequence $(R_n)_{n\in\N}$ of finite subsets of $\R$ such that $s<t$ $\Rightarrow$ $L_{R_n\cap\op]s,t\clo[}\to L_{\op]s,t\clo[}$
\end{lem}
\begin{proof} Let $((u_k,u'_k))_{k\in\N}$ be a dense sequence in
\[
\{(x,y)\in\R^2:\,x< y\}
\]
and for all $n\geq 1$, $(R_k^n)_{k\in\N}$ a nested sequence of finite subsets of $\R$ such that $\rho(L_{R_k^n\cap\op]u_k,u'_k\clo[},L_{\op]u_k,u'_k\clo[})\leqslant\sfrac1n$ for every $k\in \N$. Denote $\bigcup_{k=1}^n R^n_k$ by $R_n$. Then for any $s$ and $t>s$, $L_{R_n\cap\op]s,t\clo[}\to_{n\to\infty}L_{\op]s,t\clo[}$. Let us prove it. Fix $\eps>0$. By Lemma \ref{lem:croissance_lc} (use a subsequence $(k_n)_n$ such that $(\op]u_{k_n},u'_{k_n}\clo[)_n$ is an exhaustion of $\op]s,t\clo[$), there exists $k$ such that $\rho(L_{\op]s,t\clo[},L_{\op]u_k,u'_k\clo[})<\eps/2$. Then $L_{\op]u_k,u_k'\clo[\cap R_n}\leqlc L' \leqlc L_{\op]s,t\clo[}$, where $L'$ stands for $L_{\op]s,t\clo[\cap R_n}$ or $L_{\op]u_k,u_k'\clo[}$ for any $k$. For $n\geq\max\{k,2/\eps\}$:
\[\rho(L_{\op]s,t\clo[},L_{\op]u_k,u'_k\clo[\cap R_n})\leq \rho(L_{\op]s,t\clo[},L_{\op]u_k,u'_k\clo[})+ \rho(L_{\op]u_k,u'_k\clo[},L_{\op]u_k,u'_k\clo[\cap R_n})<\eps.\]
Therefore with $L'=L_{\op]s,t\clo[\cap R_n}$ we obtain the desired convergence to $L_{\op]s,t\clo[}$.
\end{proof}

\begin{defi}\label{defi:essentiel}Let $I$ be an interval. If, for some interval $J\supset I$ such that $J\setminus I$ is disconnected (then for all such smaller $J'$), $L_{J}\neq L_{J\setminus I}$, we call $I$ an \emph{essential atomic interval of $(\mu_t)_{t\in\R}$}. If $I=\{t\}$ is essential, we call $t$ an \emph{essential atomic time of $(\mu_t)_{t\in\R}$}.
\end{defi}

To check the parenthesis in the definition, suppose that $L_{J'}= L_{J'\setminus I}$ and get $L_{J}= L_{J\setminus I}$ by Proposition \ref{pro:compo}.

\begin{pro}\label{pro:denombrable}If a nested sequence $(R_n)_n$ is as in Lemma \ref{lem:convergence_couples} then $\bigcup_nR_n$ contains all the essential atomic times of $(\mu_t)_t$. In particular, these times are at most countably many.
\end{pro}

\begin{proof}We show a contrapositive result. Suppose that $t$ is an essential atomic time, that $s<t$ and $s'>t$ are such that $L_{\op]s,s'\clo[\setminus\{t\}}\neq L_{\op]s,s'\clo[}$, and that $(R_n)_n$ is a nested sequence of finite sets such that $t\not\in\bigcup_nR_n$, that $L_{R_n\cap\op]s,t\clo[}\to L_{\op]s,t\clo[}$ and that $L_{R_n\cap\op]t,s'\clo[}\to L_{\op]t,s'\clo[}$. Then:
\begin{align*}
L_{R_n\cap\op]s,s'\clo[}&=L_{R_n\cap\op]s,t\clo[}{}\ldot L_{R_n\cap\op]s,t\clo[}\ \ \text{by Proposition \ref{pro:compo}}\\
&\to L_{\op]s,t\clo[}{}\ldot L_{\op]t,s'\clo[}\ \ \text{by assumption and by Proposition \ref{pro:continuity}}\\
&=L_{\op]s,s'\clo[\setminus\{t\}}\neq L_{\op]s,s'\clo[}\ \ \text{by Proposition \ref{pro:compo} and the assumption}.
\end{align*}
Therefore, $(R_n)_n$ cannot be as in Lemma \ref{lem:convergence_couples}.
\end{proof}

We can now prove Theorem \ref{them:deuxieme_sec}\eqref{item:them:deuxieme_a}.

\begin{proof}[Proof of Theorem \ref{them:deuxieme_sec}\eqref{item:them:deuxieme_a}]\label{proof:themBa} Take $(R_n)_n$ a sequence as given by Lemma \ref{lem:convergence_couples}. Then, for any $s$ and $t>s$:
\begin{align*}
L_{R\cap\op]s,t\clo[}&=\lcsup\{L_{R_n\cap\op]s,t\clo[}\}\ \ \text{by Lemma \ref{lem:croissance_lc}}\\
&\geqlc\lcsup\{L_{R'_n\cap\op]s,t\clo[}\}\ \ \text{as $R'_n\subset R_n$ and by Lemma \ref{lem:existenceT}\eqref{p1:lem:existenceT}}\\
&=\lim_{n\to\infty}L_{R'_n\cap\op]s,t\clo[}=L_{\op]s,t\clo[}\ \ \text{\begin{tabular}[t]{l}by Lemma \ref{lem:croissance_lc} and\\ by property of the sets $R'_n$.\end{tabular}}
\end{align*}
But by definition, $L_{R\cap\op]s,t\clo[}\leqlc L_{\op]s,t\clo[}$, thus $L_{R\cap\op]s,t\clo[}= L_{\op]s,t\clo[}$.
\end{proof}

To provide a clear proof of Theorem \ref{them:deuxieme_sec}\eqref{item:them:deuxieme_b} we introduce an auxiliary notion in Definition \ref{defi:fits} below.
Suppose that some $P\in \ma(\mu,\mu_1,\ldots,\mu_k)$ is disintegrated as $P=\join(\mu,k_{P})$ and that $g:\R\to \R$ and $h:\R^k\to \R^k$ are measurable maps. May we disintegrate $(g\otimes h)_\#P$? In case $g$ is into one easily checks $(g\otimes \id)_\#P=\join(g_\#\mu,k^g_P)$, where $k^g_P$ is defined by $k^g_P(y,\cdot)= k^g_P(g^{-1}(y),\cdot)$, $g_\#\mu$-almost surely. Otherwise, the next notion and lemma will enable us to obtain a similar disintegration, and associated properties.

\begin{defi}\label{defi:fits}
We say that $g:\R\to \R$ \emph{fits} $P\in \ma(\mu,\mu_1,\ldots,\mu_k)$ if there exists a kernel $k^g_P$ such that $k^g_P(g(x),\cdot)=k_{P}(x,\cdot)$, $\mu$-almost surely.
\end{defi}

\skpt
\begin{rem}\label{rem:random_from_center}\begin{enumerate}
\item
\label{item:rem:random_from_center_a} If $g$ fits $P=\join(\mu,k)$ and $h$ is a measurable map from $\R^k$ into itself we can disintegrate $(g\otimes h)_\#P$ as $\join(g_\#\mu,h_\#k^g_P)$, where $h_\#k^g_P(y,\cdot)= k^g_P(y,h^{-1}(\cdot))$.

\item
\label{item:rallonge} If $g$ fits $P$, it also fits $P\ldot P'$ and $P\circ P'$.

\item
\label{item:rem:random_from_center_commute_avec_circ} If $g$ fits $P\in \ma(\mu,\nu_1,\ldots,\nu_k)$ or $Q\in\ma(\mu,\nu'_1,\ldots,\nu'_{k'})$, and if $f:\R^k\to\R^k$ and $h:\R^{k'}\to\R^{k'}$ are measurable maps, then:
\[(f\otimes g \otimes h)_\#(\trans P\circ Q)=(f\otimes g)_\#\trans P\circ (g \otimes h)_\#Q.\]
\end{enumerate}

The proofs are direct, using the definitions of a kernel, composition and concatenation in Section~\ref{subsec:caten}. Notice that, in case $\mu=\la\lfloor_{[0,1]}$ and $g$ is a quantile function (which is the only case in which we will use the remark), point \eqref{item:rem:random_from_center_commute_avec_circ} is a particular case of Lemma \ref{lem:distributivite}\eqref{item:lem:distributiviteB} below. In the language of this lemma, ``$g$ fits $\trans P$ or $Q$'' means that $\trans P\ldot L=\trans P$ or $L\ldot Q=Q$, which both imply in particular that $\trans P\ldot L\ldot Q=\trans P\ldot Q$.
\end{rem}

We will need the following little technical result.

\begin{lem}\label{lem:int_fg}Take $a<b$ in $\R$ and $f$ and $g$ two positive bounded increasing functions from $\op]a,b\clo[$ to $\R$. Then
\[
\frac1{b-a}\int_a^b fg\,\dd\la\geq \biggl(\frac1{b-a}\int_a^b f\dd\la\biggr)\biggl(\frac1{b-a}\int_a^b g\dd\la\biggr),
\]
with equality if and only if $f$ or $g$ is constant.
\end{lem}

\begin{proof} The %two
measures $\mu=f\la$ and $\nu=\frac1{b-a}\int_a^b f\,\dd\la$ have the same mass $\int_a^b f\,\dd\la$. Besides, $\nu\leqs\mu$; indeed, for any $c\in\op]a,b\clo[$, since $f$ is increasing:
\[\frac1{c-a}\mu(\op]-\infty,c])=\frac1{c-a}\int_{a}^c f\,\dd\la\leq\frac1{b-a}\int_{a}^b f\,\dd\la=\frac1{c-a}\nu(\op]-\infty,c]).\]
Hence, as $g$ is increasing, $\int g\,\dd\mu\geq\int g\,\dd\nu$, equivalent to the wished inequality.

Plainly, equality occurs if $f$ or $g$ is constant. If $f$ is not constant, we in fact proved that equality in $\int g_c \dd \mu\leq \int g_c \dd \nu$ holds for $g_c=\one_{[c,b]}$ if and only if $c\in\{a,b\}$. This remains true for positive combinations of functions $g_c$ and, being a little careful, for limits of them.
\end{proof}

\begin{lem}\label{lem:distributivite} Take $\mu\!\in\!\p(\R)$, denote by $g\!=\!G_\mu$ its quantile function and \hbox{$F\!=\!F_\mu$} its cumulative distribution function. Take $L=\como(\la,\mu)\ldot\como(\mu,\la)$ and, similarly as in Nota\-tion \ref{nota:Asx}, let $(b_i)_{i\in{\mathcal{I}}}$ be the atoms of $\mu$ and for each $i$, $A_i$ be the interval $\op]F(b_i^-),F(b_i)\clo[$ of quantile levels merged by $g$ on $b_i$. Finally set $A=\bigcup_{i\in\mathcal{I}}A_i$. Moreover, take\vspace*{-3pt}
\begin{gather*}
((\nu_i)_i,(\nu'_i)_i)\in\p(\R)^{k+k'},\\
P\in \ma(\la\lfloor_{[0,1]},\nu_1\otimes\ldots\otimes\nu_k)\AND Q\in\ma(\la_{[0,1]},\nu'_1\otimes\ldots\otimes\nu'_{k'}).
\end{gather*}
Suppose that $P$ and $Q$ have increasing kernel (for $\leqlc$ in place of $\leqs$ if $k\geq 2$ or $k'\geq2$). Then:
\begin{enumerate}
\item
\label{item:lem:distributiviteA} $\trans P\ldot L\ldot Q$ equals $\trans P\ldot Q$ if and only if for each $i\in{\mathcal I}$, at least one of the two kernel functions $x\mto k_P(x,\ \cdot\ )$ or $x\mto k_Q(x,\ \cdot\ )$ is constant on $A_i$,

\item
\label{item:lem:distributiviteB} For any measurable maps $f:\R^k\to\R^k$ and $h:\R^{k'}\to\R^{k'}$,
\[
\trans P\ldot L\ldot Q=\trans P\ldot Q\implies (f\otimes g \otimes h)_\#(\trans P\circ Q)=(f\otimes g)_\#\trans P\circ (g \otimes h)_\#Q.
\]
\end{enumerate}
\end{lem}

\begin{proof} We recall particularly here that we have throughout in mind the analogy between composition of transport plans, or of their kernels, and product of matrices, hinted at in Remark \ref{rem:matrix}. In the following, $B_1\subset\R^k$ and $B_3\subset\R^{k'}$ stand for any sets of the type $\prod_j\op]-\infty,d_j]$ and $B_2$ for any interval $\op]-\infty,b]$; $(a,b,c)$ stand for variables in the target of $(f,g,h)$, and $(x,y,z)$ for variables in their sources. The coupling $L$, equal to $\trans q_r\ldot q_r$ (with $\mu=\mu_r$) introduced in Remark \ref{rem:finesse}, is described in this remark. It is such that $\trans\!P\ldot L\ldot Q=(\trans\!P\ldot L)\ldot(L\ldot Q)=\trans((g\otimes \Id_{k})_\#P)\ldot((g\otimes \Id_{k'})_\#Q)$. In turn, for any measurable functions $f:\R^k\to\R^k$ and $h:\R^{k'}\to\R^{k'}$:\vspace*{-3pt}
\begin{equation}\label{eq:prelim_distributivite}
\begin{aligned}
\int_{B_2}&k_{(g\otimes f)_\#P}(b,B_1)\ldot k_{(g \otimes h)_\#Q}(b,B_3)\,\dd\mu(b)\\
&=\sum_{i\in\mathcal{I},\,b_i\in B_2}\mu(b_i)[k_{(g\otimes f)_\#P}(b_i,B_1)\ldot k_{(g\otimes h)_\#Q}(b_i,B_3)]\\
&\hspace*{1cm}+\int_{B_2\setminus\oldbigcup_i\{b_i\}}k_{P}(g^{-1}(b),f^{-1}(B_1))\ldot k_{Q}(g^{-1}(b),h^{-1}(B_3))\,\dd\mu(b)\\[-5pt]
&\hspace*{6.5cm}\text{as $g$ is injective outside of $A$}\\
&=\sum_{i\in \mathcal{I},\,b_i\in B_2}\frac1{\la(A_i)}\int_{A_i}k_{P}(y,f^{-1}(B_1))\,\dd\la(y)\int_{A_i}k_{Q}(y,h^{-1}(B_3))\,\dd\la(y)\\
&\hspace*{1cm}+\int_{g^{-1}(B_2)\setminus A}k_{P}(y,f^{-1}(B_1))\ldot k_{Q}(y,h^{-1}(B_3))\,\dd\la(y).
\end{aligned}
\end{equation}
Then, using for instance the expression \eqref{eq:concatenation3} given in Definition \ref{defi:concatenation} of the concatenation, and the fact that for any transport plans $R$ and $R'$, $R\ldot R'=\pr^{1,3}_\#(R\circ R')$ we get:\vspace*{-3pt}
\begin{align*}
(\trans P\ldot Q&-\trans((g\otimes \Id)_\#P)\ldot(g\otimes \Id)_\#Q)(B_1\times B_3)\\
&=\int_{[0,1]}k_{P}(y,B_1)k_Q(y,B_3)\dd y-\int_{\R}k_{(g\otimes \Id)_\#P}(b,B_1)k_{(g \otimes {\rm Id})_\#Q}(b,B_3)\dd\mu(b)\\
&=\sum_{i\in\mathcal{I}}\int_{A_i}k_{P}(y,B_1)k_{Q}(y,B_3)\dd y-\frac1{\la(A_i)}\int_{A_i}k_{P}(y,B_1)\dd y\int_{A_i}k_{Q}(y,B_3)\dd y
\end{align*}
by \eqref{eq:prelim_distributivite} with $B_2=\R$, $f=\Id$, $h=\Id$---notice in particular that the two terms of the difference, viewed as integrals on $[0,1]$, differ possibly only on $A\subset[0,1]$. Now, after Lemma \ref{lem:int_fg} applied on each $A_i$ with $f_i:y\mto k_{P}(y,B_1)$ and $g_i:y\mto k_{Q}(y,B_3)$, which are increasing since $P$ and $Q$ have increasing kernel, each term of the sum is non-negative. Thus, equality holds if and only if each vanishes, which again by Lemma~\ref{lem:int_fg} means that for each $i$, $f_i$ or $g_i$ is constant. This, holding for any $B_1$ and~$B_3$, means point \eqref{item:lem:distributiviteA}. For point \eqref{item:lem:distributiviteB}:\vspace*{-3pt}
\begin{align*}
\trans((g\otimes f)_\# P)&\circ ((g \otimes h)_\#Q)(B_1\times B_2\times B_3)\\
&=\int_{B_2} k_{(f\otimes g)_\#P}(b,B_1)k_{(g\otimes h)_\#Q}(b,B_3)\,\dd \mu(b)\quad\text{by \eqref{eq:concatenation3}}\\
&=\int_{g^{-1}(B_2)\setminus A} k_{P}(y,f^{-1}(B_1))k_{Q}(y,h^{-1}(B_3))\dd y\\
&\!\!\!\!+\sum_{i\in \mathcal{I},\,b_i\in B_2}\!\!\frac1{\la(A_i)}\int_{A_i}k_{P}(y,f^{-1}(B_1))\dd y\int_{A_i}k_{Q}(y,h^{-1}(B_3))\dd y\quad\text{by \eqref{eq:prelim_distributivite},}
\end{align*}
but on each $A_i$, by \eqref{item:lem:distributiviteA}, $k_{P}(\,.\,,f^{-1}(B_1))$ or $k_{Q}(\,.\,,h^{-1}(B_3))$ is a constant function, so~the last term may be written:\vspace*{-3pt}
\[\sum_{i\in \mathcal{I},\,b_i\in B_2}\!\int_{A_i}k_{P}(y,f^{-1}(B_1))k_{Q}(y,h^{-1}(B_3))\dd y,\]
and therefore the initial expression equals $\int_{g^{-1}(B_2)} k_{P}(y,f^{-1}(B_1))k_{Q}(y,h^{-1}(B_3))\dd y$, that is to say $(f\otimes g \otimes h)_\#(\trans P\circ Q)(B_1\times B_2\times B_3)$.
\end{proof}

Let us prove point \eqref{item:them:deuxieme_b}. Take $R\subset\R$ and $\lev_R$ and $\lev$ in $\ma((\la_t)_{t\in\R})$ given by Definition \ref{defi:lev}. We need the following claim, that extends Remark \ref{rem:mini_et_Q}, and relies eventually on Remark \ref{rem:finesse}\eqref{p2:rem:finesse}:

\begin{claim}
If $R$ is finite, $G_\#\lev_{R}=\como_{[R]}$, where $G=(\otimes_{t\in\R} G_t)$.
\end{claim}

Let us prove it. Take $R\subset\R$ finite and $n$ its cardinal. We must prove that for any finite $S\!=\!\bigcup_{i=1}^N\{s_i\}$, $G_\#\lev_R^{s_1,\ldots,s_N}\!=\!\como_{[R]}^{s_1,\ldots,s_N}$, where, by an abuse of notation we~will often make use of, $G$ stands for $\otimes_{i=1}^N G_{s_i}$. It suffices to prove it in the case $S\supset R$, which we suppose now. We introduce the cardinals $\ell_0,\ldots,\ell_n$ of the subsets of $S\setminus R$ situated between to consecutive elements $r_k$ and $r_{k+1}$ of $R\cup\{\pm\infty\}$. We re-index these subsets as $\{s_1^k,\ldots,s_{\ell_k}^k\}$. Since $\lev_R$ is Markov, $(\pr^S)_\#\lev_R=B_0\circ\cdots\circ B_m$, where:
\begin{align*}
B_k=L_{r_{k}}\circ\id_{\ell_{k}}\circ L_{r_{k+1}}=
\begin{cases}
\id_{\ell_0}\circ L_{r_1}&\text{if }k=0\\
L_{r_m}\circ \id_{\ell_m}&\text{if }k=m\\
L_{[r_k,r_{k+1}]}=L_{r_{k}}. L_{r_{k+1}}&\text{if }\ell_k=0\\
L_{[r_k,r_{k+1}]}=L_{r_{k}}\circ L_{r_{k+1}}&\text{if }\ell_k=1\\
L_{r_{k}}\circ\underbrace{\id_2\circ\cdots\circ\id_2}_{\ell_k-1}\circ L_{r_{k+1}}&\begin{nnarray}[t]{l}\text{otherwise,}\\\text{\quad \ie if }\ell_k\geq 2,\end{nnarray}
\end{cases}
\end{align*}
hence we must prove that: $G_\#(B_1\circ\cdots\circ B_m)=\como_{[R]}^{s_1,\ldots,s_N}$. For all $r\in R$, the quantile function $G_r=G_{\mu_r}$ fits $q_r=\join(\mu_r,x\mto \delta_{G_r(x)})$, so we may take $k^{G_r}_{q_r}:x\mto \delta_x$ as given by Definition \ref{defi:fits}. Now, by Remark \ref{rem:random_from_center}\eqref{item:rallonge}, it also fits $L_{\{r\}}=\trans L_{\{r\}}=q_r\ldot\trans{q_r}$. Therefore, by Remark \ref{rem:random_from_center}\eqref{item:rallonge} and \eqref{item:rem:random_from_center_commute_avec_circ}, $G_\#(B_1\circ\cdots\circ B_m)=B'_1\circ\ldots\circ B'_m$, where:
\begin{align*}
B'_k&=G_\# B_k=\begin{cases}
(\otimes_{s\leqslant r_1} G_{s})_\#\id_{\ell_0}\circ L_{r_1}&\text{if }k=0\\
(\otimes_{s\geqslant r_m} G_{s})_\#L_{r_m}\circ \id_{\ell_m}&\text{if }k=m\\
(G_{r_k}\otimes G_{r_{k+1}})_\#(L_{r_{k}}{}\ldot L_{r_{k+1}})&\text{if }\ell_k=0\\
(G_{r_k}\otimes G_{s_1^k}\otimes G_{r_{k+1}})_\#(L_{r_{k}}\circ L_{r_{k+1}})&\text{if }\ell_k=1\\
(\otimes_{r_k\leqslant s\leqslant r_{k+1}}G_s)_\#(L_{r_{k}}\circ\id_{\ell_k}\circ L_{r_{k+1}})&\text{otherwise, \ie }\ell_k\geq 2.
\end{cases}
\end{align*}

\noindent Proving that $B'_k=\como(\mu_{r_k},\mu_{s^k_1},\ldots,\mu_{s^k_{\ell_k-1}},\mu_{r_{k+1}})$ will now prove the claim. For simplicity we assume $k\notin \{0,m\}$ and $\ell_k\geq 1$ but the other cases, which are simpler, can be proved similarly. Observe that:
\begin{itemize}
\item
by definition of $\como$ and $G$,
\[
\como(\mu_{r_k},\mu_{s_1^k},\ldots,\mu_{s^k_{\ell_k}},\mu_{r_{k+1}})=\law((G_{r_k}(U),\ldots,G_{r_{k+1}}(U)),
\]
where $U$ is a variable of law $\la$ on $[0,1]$,

\item
$(\otimes_{r_k\leqslant s\leqslant r_{k+1}}G_s)_\#(L_{r_{k}}\circ\id_{\ell_k}\circ L_{r_{k+1}})$ is the law of
\[
(G_{r_k}(U_1),G_{s^k_1}(U_2),\ldots,G_{s^k_{\ell_k}}(U_2),G_{r_{k+1}}(U_3)),
\]
where $\law(U_1,U_2,\ldots,U_2,U_3)=B_k$. In particular, $\law(U_i)=\la$ for $i=1,2,3$, $\law(U_1,U_2)=L_{r_k}$, and $\law(U_2,U_3)=L_{r_{k+1}}$.
\end{itemize}
So only the first and last variables may differ. Now, notice that by Remark \ref{rem:finesse}\eqref{p2:rem:finesse} applied to $T=L_{\{r\}}=q_r\ldot\trans q_r$, if two variables $(U,V)$ satisfy $\law(U,V)=L_{\{r\}}$, the law of $(G_{r}(U),G_r(V))$ is
\[
\trans q_r\ldot L_{\{r\}}{}\ldot q_r=\trans q_r\ldot q_r\ldot\trans q_r\ldot q_r=\trans q_r\ldot q_r=\id_{\mu_r,2},
\]
therefore $G_{r}(U)=G_r(V)$ almost surely. Using this with $r={r_k}$ and $(U,V)=(U_1,U_2)$, respectively $r=r_{k+1}$ and $(U,V)=(U_2,U_3)$, we get respectively $G_{r_k}(U_1)=G_{r_k}(U_2)$ and $G_{r_{k+1}}(U_2)=G_{r_{k+1}}(U_3)$ almost surely. The claim is proved.

Now take $R\subset \R$ satisfying \eqref{eq:convergence_doug} and $(R_n)_n$ a nested exhaustion of it by finite sets. Let $\widetilde\lev\in\ma((\la_t)_{t\in\R})$ be the Markov process having $\widetilde\lev{}^{s,t}= L_{R\cap[s,t]}$ as 2\nobreakdash-mar\-ginals. Recall that $L_{R\cap[s,t]}=L_{R\cap \{s\}}{}\ldot L_{R\cap\op]s,t\clo[}{}\ldot L_{R\cap \{t\}}=L_{R\cap \{s\}}{}\ldot L_{\op]s,t\clo[}{}\ldot L_{R\cap \{t\}}=\lim_{n\to\infty}\lev_{R_n\cap[s,t]}=\lim_{n\to\infty}\lev_{R_n}^{s,t}$. It exists by Corollary \ref{cor:consist}, whose assumption is satisfied by Proposition \ref{pro:compo}. All these transport plans
have increasing kernel, hence by Lemma \ref{lem:closed}\label{appel_lemme_closed}, for every $S=\bigcup_{i=1}^k\{s_i\}$:
\begin{multline}
\lev_{R_n}^{s_1,\ldots,s_k}=\lev_{R_n}^{s_1,s_2}\circ\cdots\circ \lev_{R_n}^{s_{k-1},s_k}\label{eq:lev_egal_concatenation}\\
\to\widetilde \lev^{s_1,s_2}\circ\cdots\circ \widetilde\lev{}^{s_{k-1},s_k}=\widetilde\lev{}^{s_1,\ldots,s_k}.
\end{multline}
Thus $\lev_{R_n}$ converges weakly to $\widetilde\lev$, and then by Remark \ref{rem:diese_contractant}, $G_\# \lev_{R_n}\to G_\# \widetilde\lev$. We are left with the tasks to prove $G_\# \lev_{R_n}=\como_{[R_n]}$ and $G_\# \widetilde\lev=\mq$. The former is our claim above. Let us prove the latter.

\begin{note}
\label{lieu_preuve_glev} At the beginning of Section~\ref{sec:construction} we announced \eqref{eq:Glev}, \ie $G_\#\lev=\mq$. In~fact, we prove $G_\# \widetilde\lev=\mq$, which is a bit more difficult. To get \eqref{eq:Glev} the same arguments work, the final reasoning with Lemma \ref{lem:distributivite}\eqref{item:lem:distributiviteB} being replaced by a direct use of Remark \ref{rem:random_from_center}\eqref{item:rem:random_from_center_commute_avec_circ}, as for all $\{r_1,\ldots,r_\ell\}$, $G_{r_\ell}$ fits $\lev^{r_1,\ldots,r_\ell}$.
\end{note}

We recall that $\mq$ was defined as the unique Markov law with the same marginals of dimension 2 as $G_\# \lev$. But by \eqref{eq:convergence_doug} and Proposition \ref{pro:comp_conv}\eqref{ccc}, the 2-marginals of $G_\#\widetilde\lev$ and $G_\#\lev$ are equal. Hence it is sufficient to prove that
$G_\# \widetilde \lev$ is Markov, \ie that for all $(s_1,\ldots,s_k)$, $(G_\#\widetilde\lev)^{s_1,\ldots,s_k}=(G_\#\widetilde\lev)^{s_1,s_2}\circ\ldots\circ(G_\#\widetilde\lev)^{s_{1-1},s_k}$. Since $\widetilde\lev$ is Markov, $(\widetilde\lev)^{s_{2},\ldots,s_{k}}=\widetilde\lev{}^{s_{2},s_{3}}\circ\ldots\circ\widetilde\lev{}^{s_{k-1},s_{k}}$; besides, notice the following fact, that we will prove a bit below:

\begin{fact}
For all $(s_i)_{i=1}^k\in\R^k$, $\lev^{s_1,\ldots,s_k}$ viewed as a transport plan from $(\R,\la)$ to $(\R^{k-1},\la^{\otimes k-1})$ has increasing kernel (for the order $\leqlc$ instead of $\leqs$).
\end{fact}

Hence we may conclude by using $k-1$ times Lemma \ref{lem:distributivite}\eqref{item:lem:distributiviteB}. Let us check the first step. Since the measures $\trans(\widetilde\lev{}^{s_1,s_2})$ and $\widetilde\lev{}^{s_2,\ldots,s_k}$ have increasing kernel, we have only to show that $\widetilde\lev{}^{s_1,s_2}.\widetilde\lev{}^{s_{2},s_3}\!=\!\widetilde{\lev}^{s_1,s_{2}}{}\ldot L_{\{s_{2}\}}.\widetilde{\lev}^{s_{2},s_3}$, \ie by definition of~$\widetilde\lev$, that $L_{R\cap[s_1,s_2]}{}\ldot L_{R\cap[s_2,s_3]}=L_{R\cap[s_1,s_2]}{}\ldot L_{s_2}{}\ldot L_{R\cap[s_2,s_3]}$. This amounts to checking that:\vspace*{-3pt}
\[
\begin{cases}
L_{\semi s_1,s_2\clo[}{}\ldot L_{\op]s_2,s_3\semi}=L_{\semi s_1,s_2\clo[}{}\ldot L_{s_2}{}\ldot L_{\op]s_2,s_3\semi}&\ \text{if}\ s_2\not\in R\\ L_{\semi s_1,s_2]}{}\ldot L_{[s_2,s_3\semi}=L_{\semi s_1,s_2]}{}\ldot L_{s_2}{}\ldot L_{[s_2,s_3\semi}&\ \text{otherwise.}\end{cases}
\]
The second point is true by Proposition \ref{pro:compo}, and if $s_2\not\in R$, by Proposition \ref{pro:denombrable}, $s_2$ is not an essential atomic time, which is the wanted equality. We finally must prove the fact stated above. Actually it is true for any concatenation $P_1\circ\ldots\circ P_k$ of couplings $P_i$ from $\R$ to $\R$ with increasing kernel —and $\lev^{s_1,\ldots,s_k}$ is of this type, see~\eqref{eq:lev_egal_concatenation}. We check this for $k=2$; the same argument, applied by induction, gives the general case. Take $B_2$ and $B_3$ two intervals of the type $\op]-\infty,a]$ and $x\leqslant x'$; we must show that: $k_{P_1\circ P_2}(x,B_2\times B_3)\geqslant k_{P_1\circ P_2}(x',B_2\times B_3)$. As $P_2$ has increasing kernel, the function $k_{P_2}(\espcdot,B_3)$ is decreasing (thus also $\one_{B_2}k_{P_2}(\espcdot,B_3)$). As $P_1$ has increasing kernel, $k_{P_1}(x,\espcdot)\leqs k_{P_1}(x',\espcdot)$. Thus:\vspace*{-3pt}
\[k_{P_1\circ P_2}(x,B_2\times B_3)\!=\!\int k_{P_1}(x,\dd y) \one_{B_2}k_{P_2}(y,B_3)\!\geqslant\! \int k_{P_1}(x',\dd y) \one_{B_2}k_{P_2}(y,B_3),\]
the desired result. This proves part \eqref{item:them:deuxieme_b} of the theorem.

\section{Construction and characterization of the Markov-quantile process}\label{sec:construction}

\emph{Now $(\mu_t)_{t\in \R}$ is a family of probability measures on $\R$, $F_t$ denotes the cumulative distribution function $F_{\mu_t}$ of $\mu_t$ and $G_t$ its quantile function, see Section~\ref{ss:minimality}.} In this section we build the Markov-quantile measure $\mq$ and prove Theorems A and B. Our proof of Theorems A--\ref{them:b} is based on transport plans $L_{[s,t]}\in \ma(\la,\la)$, defined in Section~\ref{subsec:little_t}, that are the 2-marginals of an important auxiliary process law called the quantile level measure $\lev\in\ma(\la_t)_{t\in\R}$, where each $\lambda_t$ is a copy indexed by $t$ of $\la=\la_{[0,1]}$. Here is the link between $\lev$ and $\mq$: $G_t$ maps $([0,1],\la_t)$ to $(\R\cup\{\pm \infty\},\mu_t)$; set $G=\otimes_{t\in\R}G_t$, so that $G_\#$ maps $\ma((\lambda_t)_{t\in\R})$ to $\ma((\mu_t)_{t\in\R})$. Then, as proved in the proof Theorem \ref{them:deuxieme_sec} p.\,\pageref{lieu_preuve_glev}:\vspace*{-3pt}
\begin{equation}\label{eq:Glev}
G_\#\lev=\mq.
\end{equation}

\begin{thebibliography}{99}
\bibitem{Albin} J. M. P. Albin - “A continuous non-Brownian motion martingale with Brownian motion marginal distributions”, Statist. Probab. Lett. 78 (2008) no. 6, p. 682-686
\bibitem{AGS} L. Ambrosio, N. Gigli \& G. Savaré - Gradient flows in metric spaces and in the space of probability measures, Lectures in Math. ETH Zürich, Birkhäuser Verlag, Basel, 2008
\bibitem{BHS} M. Beiglböck, M. Huesmann \& F. Stebegg - “Root to Kellerer”, in Séminaire de Probabilités XLVIII, Lect. Notes in Math., vol. 2168, Springer, Cham, 2016, p. 1-12
\bibitem{BJ} M. Beiglböck \& N. Juillet - “On a problem of optimal transport under marginal martingale constraints”, Ann. Probab. 44 (2016) no. 1, p. 42-106
\bibitem{BeJu16} M. Beiglböck \& N. Juillet - “Shadow couplings”, Trans. Amer. Math. Soc. 374 (2021) no. 7, p. 4973-5002
\bibitem{BeBr} J.-D. Benamou \& Y. Brenier - “A numerical method for the optimal time-continuous mass transport problem and related problems”, in Monge Ampère equation: applications to geometry and optimization (Deerfield Beach, FL, 1997), Contemp. Math., vol. 226, American Mathematical Society, Providence, RI, 1999, p. 1-11
\bibitem{Bi} P. Billingsley - Convergence of probability measures, Wiley Series in Probability and Statistics, John Wiley \& Sons Inc., New York, 1999
\bibitem{BouJui2} C. Boubel \& N. Juillet - “On absolutely continuous curves in the Wasserstein space over \ensuremath{\mathbb{R}} and their representation by an optimal Markov process”, 2021
\bibitem{BrSh} G. Brunick \& S. Shreve - “Mimicking an Itô process by a solution of a stochastic differential equation”, Ann. Appl. Probab. 23 (2013) no. 4, p. 1584-1628
\bibitem{Doe} W. Doeblin - “Sur certains mouvements aléatoires discontinus”, Skand. Aktuarietidskr. 22 (1939), p. 211-222
\bibitem{Dupire} B. Dupire et al. - “Pricing with a smile”, Risk 7 (1994) no. 1, p. 18-20
\bibitem{Fe} H. Federer - Geometric measure theory, Grundlehren Math. Wiss., vol. 153, Springer-Verlag New York Inc., New York, 1969
\bibitem{Fe_t1} W. Feller - An introduction to probability theory and its applications. Vol. I, John Wiley \& Sons, Inc., New York-London-Sydney, 1968
\bibitem{Gy} I. Gyöngy - “Mimicking the one-dimensional marginal distributions of processes having an Itô differential”, Probab. Theory Relat. Fields 71 (1986) no. 4, p. 501-516
\bibitem{HaKl} K. Hamza \& F. C. Klebaner - “A family of non-Gaussian martingales with Gaussian marginals”, J. Appl. Math. Stochastic Anal. (2007), article ID 92723, 19 pages
\bibitem{HTT} P. Henry-Labordère, X. Tan \& N. Touzi - “An explicit martingale version of the one-dimensional Brenier’s theorem with full marginals constraint”, Stochastic Process. Appl. 126 (2016) no. 9, p. 2800-2834
\bibitem{Hil_ejp} E. Hillion - “W 1,+ -interpolation of probability measures on graphs”, Electron. J. Probab. 19 (2014), article ID 92, 29 pages
\bibitem{HPRY} F. Hirsch, C. Profeta, B. Roynette \& M. Yor - Peacocks and associated martingales, with explicit constructions, Bocconi \& Springer Series, vol. 3, Springer, Milan, 2011
\bibitem{HiRo13} F. Hirsch \& B. Roynette - “On \ensuremath{\mathbb{R}} d -valued peacocks”, ESAIM Probab. Statist. 17 (2013), p. 444-454
\bibitem{HiRoYo14} F. Hirsch, B. Roynette \& M. Yor - “Kellerer’s theorem revisited”, in Asymptotic Laws and Methods in Stochastics. Volume in Honour of Miklos Csorgo, Fields Inst. Commun. Series, vol. 76, Fields Inst. Res. Math. Sci., Toronto, ON, 2015, p. 347-363
\bibitem{Hob_fake} D. G. Hobson - “Fake exponential Brownian motion”, Statist. Probab. Lett. 83 (2013) no. 10, p. 2386-2390
\bibitem{Hob16} D. G. Hobson - “Mimicking martingales”, Ann. Appl. Probab. 26 (2016) no. 4, p. 2273-2303
\bibitem{Ju_seminaire} N. Juillet - “Peacocks parametrised by a partially ordered set”, in Séminaire de Probabilités XLVIII, Lect. Notes in Math., vol. 2168, Springer, Cham, 2016, p. 13-32
\bibitem{Ju_ejp} N. Juillet - “Martingales associated to peacocks using the curtain coupling”, Electron. J. Probab. 23 (2018), article ID 8, 29 pages
\bibitem{Kal} O. Kallenberg - Foundations of modern probability, Probability and its Appl. (New York), Springer-Verlag, New York, 2002
\bibitem{KaKr} T. Kamae \& U. Krengel - “Stochastic partial ordering”, Ann. Probab. 6 (1978) no. 6, p. 1044-1049
\bibitem{KLN} K. S. Kaminsky, E. M. Luks \& P. I. Nelson - “Strategy, nontransitive dominance and the exponential distribution”, Austral. J. Statist. 26 (1984) no. 2, p. 111-118
\bibitem{Ke72} H. G. Kellerer - “Markov-Komposition und eine Anwendung auf Martingale”, Math. Ann. 198 (1972), p. 99-122
\bibitem{Ke73} H. G. Kellerer - “Integraldarstellung von Dilationen”, in Transactions of the Sixth Prague Conference on Information Theory, Statistical Decision Functions, Random Processes (Tech. Univ., Prague, 1971; dedicated to the memory of Antonín Špaček), Academia, Prague, 1973, p. 341-374
\bibitem{Ke86} H. G. Kellerer - “Order conditioned independence of real random variables”, Math. Ann. 273 (1986), p. 507-528
\bibitem{Ke87} H. G. Kellerer - “Markov property of point processes”, Probab. Theory Relat. Fields 76 (1987), p. 71-80
\bibitem{Leo16} C. Léonard - “Lazy random walks and optimal transport on graphs”, Ann. Probab. 44 (2016) no. 3, p. 1864-1915
\bibitem{Lis} S. Lisini - “Characterization of absolutely continuous curves in Wasserstein spaces”, Calc. Var. Partial Differential Equations 28 (2007) no. 1, p. 85-120
\bibitem{Low2} G. Lowther - “Fitting martingales to given marginals”, 2008
\bibitem{Low} G. Lowther - “Limits of one-dimensional diffusions”, Ann. Probab. 37 (2009) no. 1, p. 78-106
\bibitem{MaYo02} D. B. Madan \& M. Yor - “Making Markov martingales meet marginals: with explicit constructions”, Bernoulli 8 (2002) no. 4, p. 509-536
\bibitem{Naga} M. Nagasawa - Schrödinger equations and diffusion theory, Monographs in Math., vol. 86, Birkhäuser Verlag, Basel, 1993
\bibitem{Ol} Y. Ollivier - “Ricci curvature of Markov chains on metric spaces”, J. Funct. Anal. 256 (2009) no. 3, p. 810-864
\bibitem{pages} G. Pagès - “Functional co-monotony of processes with applications to peacocks and barrier options”, in Séminaire de Probabilités XLV, Lect. Notes in Math., vol. 2078, Springer, Cham, 2013, p. 365-400
\bibitem{Pass} B. Pass - “On a class of optimal transportation problems with infinitely many marginals”, SIAM J. Appl. Math. 45 (2013) no. 4, p. 2557-2575
\bibitem{RR} S. T. Rachev \& L. Rüschendorf - Mass transportation problems. Vol. I \& II, Probability and its Appl. (New York), Springer-Verlag, New York, 1998
\bibitem{RSY} Y. Rinott, M. Scarsini \& Y. Yu - “A Colonel Blotto gladiator game”, Math. Oper. Res. 37 (2012) no. 4, p. 574-590
\bibitem{SS} M. Shaked \& J. G. Shanthikumar - Stochastic orders, Springer Series in Statistics, Springer, New York, 2007
\bibitem{St65} V. Strassen - “The existence of probability measures with given marginals”, Ann. Math. Statist. 36 (1965), p. 423-439
\bibitem{Vi1} C. Villani - Topics in optimal transportation, Graduate Studies in Math., vol. 58, American Mathematical Society, Providence, RI, 2003
\bibitem{Vi2} C. Villani - Optimal transport, Grundlehren Math. Wiss., vol. 338, Springer-Verlag, 2009
\end{thebibliography}
\end{document}
